%==============================================================================
% MASTER THESIS — AIVANCITY PGE5
% Design, Development, and Evaluation of a RAG Conversational Assistant
% Based on the Model Context Protocol (MCP) for Secure and Intelligent
% Access to Enterprise Knowledge Stored in SharePoint
%
% Author  : Jean Direl NZE
% Program : Programme Grande École (PGE5)
% School  : Aivancity School for Technology and Business, Villejuif, France
% Year    : 2025–2026
%==============================================================================

\documentclass[12pt, a4paper, oneside]{report}

%------------------------------------------------------------------------------
% PACKAGES
%------------------------------------------------------------------------------
\usepackage[T1]{fontenc}
\usepackage[utf8]{inputenc}
\usepackage[english]{babel}
\usepackage{lmodern}
\usepackage{microtype}

% Page geometry
\usepackage[
  top=2.5cm, bottom=2.5cm,
  left=3.0cm, right=2.5cm
]{geometry}

% Typography
\usepackage{setspace}
\onehalfspacing
\usepackage{parskip}
\setlength{\parskip}{6pt}
\setlength{\parindent}{0pt}

% Mathematics
\usepackage{amsmath, amssymb, amsthm}

% Graphics and figures
\usepackage{graphicx}
\usepackage{float}
\usepackage{caption}
\usepackage{subcaption}
\usepackage[table,xcdraw]{xcolor}
\graphicspath{{figures/}}

% Tables
\usepackage{booktabs}
\usepackage{longtable}
\usepackage{array}
\usepackage{multirow}
\usepackage{tabularx}

% Code listings
\usepackage{listings}
\usepackage{minted}
\lstset{
  basicstyle=\ttfamily\footnotesize,
  breaklines=true,
  frame=single,
  captionpos=b,
  numbers=left,
  numberstyle=\tiny,
  tabsize=2
}

% Hyperlinks and PDF metadata
\usepackage[
  colorlinks=true,
  linkcolor=black,
  citecolor=black,
  urlcolor=blue,
  pdftitle={RAG Conversational Assistant Based on MCP for Enterprise SharePoint Knowledge Access},
  pdfauthor={Jean Direl NZE},
  pdfsubject={Master Thesis — Aivancity PGE5},
  pdfkeywords={RAG, MCP, Model Context Protocol, SharePoint, Enterprise AI, LLM, FAISS}
]{hyperref}

% Bibliography
\usepackage[
  backend=biber,
  style=authoryear,
  sorting=nyt,
  maxnames=3,
  minnames=1,
  natbib=true
]{biblatex}
\addbibresource{bibliography/references.bib}

% Header and footer
\usepackage{fancyhdr}
\pagestyle{fancy}
\fancyhf{}
\fancyhead[L]{\small\leftmark}
\fancyhead[R]{\small Jean Direl NZE — Aivancity PGE5}
\fancyfoot[C]{\thepage}
\renewcommand{\headrulewidth}{0.4pt}

% Chapter styling
\usepackage{titlesec}
\titleformat{\chapter}[display]
  {\normalfont\Large\bfseries}
  {\chaptertitlename\ \thechapter}{14pt}{\huge}
\titlespacing*{\chapter}{0pt}{0pt}{20pt}

% Acronyms and glossary
\usepackage[acronym,toc]{glossaries}
\makeglossaries
%==============================================================================
% GLOSSARY AND ACRONYMS
%==============================================================================

\newacronym{rag}{RAG}{Retrieval-Augmented Generation}
\newacronym{llm}{LLM}{Large Language Model}
\newacronym{mcp}{MCP}{Model Context Protocol}
\newacronym{faiss}{FAISS}{Facebook AI Similarity Search}
\newacronym{acl}{ACL}{Access Control List}
\newacronym{sso}{SSO}{Single Sign-On}
\newacronym{api}{API}{Application Programming Interface}
\newacronym{nlp}{NLP}{Natural Language Processing}
\newacronym{ui}{UI}{User Interface}
\newacronym{dnsi}{DNSI}{Direction du Numérique et des Systèmes d'Information}
\newacronym{cerp}{CERP}{Collaborative Enterprise Resource Platform}
\newacronym{dsr}{DSR}{Design Science Research}
\newacronym{pdf}{PDF}{Portable Document Format}
\newacronym{docx}{DOCX}{Microsoft Word Document Format}
\newacronym{pge}{PGE}{Programme Grande École}
\newacronym{mrr}{MRR}{Mean Reciprocal Rank}
\newacronym{ndcg}{nDCG}{Normalized Discounted Cumulative Gain}
\newacronym{sus}{SUS}{System Usability Scale}
\newacronym{tam}{TAM}{Technology Acceptance Model}
\newacronym{ann}{ANN}{Approximate Nearest Neighbor}
\newacronym{bert}{BERT}{Bidirectional Encoder Representations from Transformers}
\newacronym{gpu}{GPU}{Graphics Processing Unit}
\newacronym{cpu}{CPU}{Central Processing Unit}
\newacronym{saml}{SAML}{Security Assertion Markup Language}
\newacronym{oauth}{OAuth}{Open Authorization}
\newacronym{qa}{QA}{Question Answering}
\newacronym{kms}{KMS}{Knowledge Management System}
\newacronym{ai}{AI}{Artificial Intelligence}
\newacronym{ml}{ML}{Machine Learning}


% Theorem environments
\newtheorem{definition}{Definition}[chapter]
\newtheorem{proposition}{Proposition}[chapter]
\newtheorem{remark}{Remark}[chapter]

% Custom colors
\definecolor{aivancityblue}{RGB}{0, 82, 165}
\definecolor{lightgray}{RGB}{245, 245, 245}

%------------------------------------------------------------------------------
% DOCUMENT BEGIN
%------------------------------------------------------------------------------
\begin{document}

%---- FRONT MATTER ------------------------------------------------------------
\pagenumbering{roman}

%==============================================================================
% COVER PAGE
%==============================================================================
\begin{titlepage}
\thispagestyle{empty}

\begin{center}

% School header
{\large \textbf{aivancity School for Technology and Business}}\\[0.2cm]
{\normalsize La Grande École de l'IA \& de la Data}\\[0.1cm]
{\small Paris Île-de-France · Nice Côte d'Azur}\\[1.5cm]

\hrule height 1pt
\vspace{0.3cm}
{\small \textsc{Programme Grande École — PGE5}}\\
{\small \textsc{Final Year Project (PFE) — 2025–2026}}
\vspace{0.3cm}
\hrule height 1pt

\vspace{2cm}

% Title
{\LARGE \textbf{Design, Development, and Evaluation of a}}\\[0.4cm]
{\LARGE \textbf{Retrieval-Augmented Generation (RAG)}}\\[0.4cm]
{\LARGE \textbf{Conversational Assistant Based on the}}\\[0.4cm]
{\LARGE \textbf{Model Context Protocol (MCP) for Secure}}\\[0.4cm]
{\LARGE \textbf{and Intelligent Access to Enterprise}}\\[0.4cm]
{\LARGE \textbf{Knowledge Stored in SharePoint}}

\vspace{2cm}

% Host company
{\normalsize \textit{Host Organization:}}\\[0.2cm]
{\large \textbf{DNSI — Direction du Numérique et des Systèmes d'Information}}

\vspace{2cm}

% Author and supervisor
\begin{tabular}{ll}
  \textbf{Author:}               & Jean Direl NZE \\[0.3cm]
  \textbf{Supervising Professor:} & [Supervisor Name] \\[0.3cm]
  \textbf{Year of Submission:}    & 2026 \\
\end{tabular}

\vfill

{\small \textit{This document is submitted in partial fulfillment of the requirements}}\\
{\small \textit{for the Programme Grande École (PGE) diploma at Aivancity.}}

\end{center}
\end{titlepage}
\cleardoublepage

%==============================================================================
% ABSTRACT
%==============================================================================
\chapter*{Abstract}
\addcontentsline{toc}{chapter}{Abstract}

\textbf{Keywords:} Retrieval-Augmented Generation, Model Context Protocol, SharePoint, Enterprise Knowledge Management, Large Language Models, Semantic Search, FAISS, Security, Hallucination.

\bigskip

Large organizations accumulate substantial volumes of technical and business documentation on enterprise platforms such as SharePoint, yet access to this knowledge remains fragmented and poorly suited to complex contextual queries. Large Language Models (LLMs) offer a natural language interface to enterprise knowledge, but their direct deployment is impeded by three structural limitations: the generation of unverified facts (hallucination), the absence of any mechanism for accessing private organizational documentation, and the impossibility of enforcing document-level access control.

This PFE report addresses the following research question: \textit{to what extent does a Retrieval-Augmented Generation (RAG) architecture grounded in the Model Context Protocol (MCP) enable more reliable, secure, and traceable access to enterprise knowledge stored in SharePoint, compared to a standalone LLM and a naive RAG implementation?}

We design, implement, and evaluate AntiGravity --- a modular RAG conversational assistant that leverages MCP as a standardized protocol layer between the AI reasoning engine and the SharePoint document corpus of the DNSI (Direction du Num\'{e}rique et des Syst\`{e}mes d'Information). The architecture separates a security-aware MCP Context Server, which manages a FAISS semantic index with ACL-filtered document access, from a conversational client that combines query classification, dense retrieval, cross-encoder re-ranking, and natural language generation via a locally deployed Mistral 7B Instruct model.

The evaluation follows a four-dimensional protocol: (1) retrieval quality (Precision@K, Recall@K, MRR, nDCG) on a domain-specific 50-item benchmark constructed from the DNSI corpus; (2) generation quality (faithfulness, groundedness, hallucination rate via RAGAS); (3) system performance (end-to-end latency, component-level profiling); and (4) user acceptance (SUS, TAM, task-based evaluation with DNSI employees).
% [NOTE FOR AUTHOR: Insert the actual quantitative results here once experiments are complete.
%  Replace this comment block with measured values:
%  - Hallucination rate: LLM alone (X%) vs. Naive RAG (Y%) vs. RAG+MCP (Z%)
%  - RAGAS Faithfulness: Cond. A / Cond. B / Cond. C
%  - MCP protocol overhead at P95 (ms)
%  - End-to-end latency P95 on CPU (s) and expected GPU estimate
%  - SUS score (mean, 95% CI), N participants
%  - TAM Perceived Usefulness vs. SharePoint baseline (mean, p-value)
%  - Task completion rate and mean time saving]

The primary scientific contribution is a formal comparative evaluation framework for RAG+MCP architectures applied to a private enterprise SharePoint corpus with ACL enforcement. The engineering contribution is a reference architecture for secure, permission-aware conversational AI in enterprise environments, deployable via Docker under corporate proxy constraints.

\bigskip

\noindent\rule{\textwidth}{0.4pt}

\bigskip

\textbf{R\'{e}sum\'{e} (French)}

\textbf{Mots-cl\'{e}s~:} G\'{e}n\'{e}ration Augment\'{e}e par R\'{e}cup\'{e}ration, Model Context Protocol, SharePoint, Gestion des Connaissances en Entreprise, Grands Mod\`{e}les de Langage, Recherche S\'{e}mantique, S\'{e}curit\'{e}, Hallucination.

\bigskip

Les grandes organisations accumulent de vastes quantit\'{e}s de documentation technique et m\'{e}tier sur des plateformes comme SharePoint, mais l'acc\`{e}s \`{a} ce patrimoine informationnel reste fragment\'{e} et insuffisant pour r\'{e}pondre \`{a} des requ\^{e}tes complexes et contextuelles. Bien que les grands mod\`{e}les de langage (LLMs) offrent une interface en langage naturel prometteuse, leur d\'{e}ploiement direct est entrav\'{e} par les ph\'{e}nom\`{e}nes d'hallucination, le manque de connaissance des donn\'{e}es priv\'{e}es, et l'absence de m\'{e}canismes natifs de contr\^{o}le d'acc\`{e}s et de tra\c{c}abilit\'{e} des sources.

Cette th\`{e}se con\c{c}oit, impl\'{e}mente et \'{e}value AntiGravity --- un assistant conversationnel RAG modulaire qui exploite le Model Context Protocol (MCP) comme couche de protocole standardis\'{e}e entre le moteur d'inf\'{e}rence IA et le corpus documentaire SharePoint de la DNSI. L'architecture s\'{e}pare un serveur de contexte MCP (index FAISS s\'{e}mantique avec filtrage ACL) d'un client conversationnel combinant classification de requ\^{e}te, recherche dense, re-classement par encodeur crois\'{e} et g\'{e}n\'{e}ration via Mistral 7B. L'\'{e}valuation suit un protocole en quatre dimensions (qualit\'{e} de r\'{e}cup\'{e}ration, qualit\'{e} de g\'{e}n\'{e}ration, performance syst\`{e}me, acceptation utilisateur)~; les r\'{e}sultats quantitatifs d\'{e}taill\'{e}s sont report\'{e}s au chapitre~\ref{chap:evaluation}.

% [NOTE FOR AUTHOR: Add key figures from the evaluation here once experiments are completed.]

\bigskip
\noindent\rule{\textwidth}{0.4pt}

%==============================================================================
% ACKNOWLEDGEMENTS AND DECLARATION
%==============================================================================
\chapter*{Acknowledgements and Declaration}
\addcontentsline{toc}{chapter}{Acknowledgements and Declaration}

\section*{Acknowledgements}

[To be completed by the author upon thesis finalization.]

I wish to express my sincere gratitude to [Supervising Professor Name] for their methodological guidance and intellectual support throughout this research project.

I am deeply thankful to the DNSI team for providing access to the SharePoint corpus, for their availability during the user study, and for the trust they placed in this research.

I also thank the faculty and staff of Aivancity School for Technology and Business for the rigorous academic environment that made this work possible.

Finally, [personal acknowledgements to be completed by the author].

\bigskip

\section*{Declaration}

I, Jean Direl NZE, declare that:

\begin{enumerate}
  \item This thesis is entirely my own work, except where explicitly stated otherwise through appropriate references and citations.
  \item I have read and understood the plagiarism policy of Aivancity School for Technology and Business, and I confirm that this work does not violate that policy.
  \item I have obtained necessary permissions for any third-party material included in this work.
  \item Generative AI tools (specifically Claude, developed by Anthropic) were used as an assistive tool during the writing process. All uses are documented in Annex~E in compliance with Aivancity's policy on generative AI use. All analysis, reasoning, interpretation, and critical contributions are my own intellectual work.
  \item The DNSI organizational context, documents, and user study data are treated in accordance with the confidentiality agreement signed between myself, the DNSI, and Aivancity School for Technology and Business.
\end{enumerate}

\bigskip

\noindent Signed: \underline{\hspace{5cm}} \hspace{2cm} Date: \underline{\hspace{3cm}}

\noindent Jean Direl NZE


\tableofcontents
\listoffigures
\listoftables
\printglossary[type=\acronymtype, title=List of Abbreviations]

%---- BODY --------------------------------------------------------------------
\cleardoublepage
\pagenumbering{arabic}

%==============================================================================
% CHAPTER 1 — INTRODUCTION
%==============================================================================
\chapter{Introduction}
\label{chap:introduction}

%------------------------------------------------------------------------------
\section{General Context: The Enterprise Knowledge Access Problem}
\label{sec:intro-context}
%------------------------------------------------------------------------------

The digitalization of organizational processes has generated an unprecedented volume of business documentation. Technical specifications, procedural guides, project reports, regulatory frameworks, meeting minutes, and operational manuals accumulate continuously on enterprise collaborative platforms. This documentation represents an organization's collective intelligence — the crystallization of years of expertise, decision history, and operational knowledge. Yet despite the theoretical availability of this knowledge, organizations consistently struggle to translate \textit{stored} information into \textit{accessible} information.

The symptoms of this gap are well documented in the enterprise knowledge management literature \citep{Alavi2001, Davenport1998}. Employees spend a disproportionate fraction of their working time searching for information rather than using it. Complex queries — those requiring synthesis across multiple documents, historical context, or domain-specific reasoning — are particularly poorly served by conventional search interfaces. The result is what \citet{Davenport1998} term \textit{knowledge friction}: the resistance that organizational structures, technical interfaces, and information architectures impose between the knowledge that exists and the knowledge that can be practically applied.

Microsoft SharePoint, the enterprise content management platform used by over 200 million users globally \citep{Microsoft2024}, is a paradigmatic example of this tension. SharePoint provides robust document storage, versioning, and access control capabilities, but its search functionality is fundamentally lexical: it matches keywords without understanding intent, context, or semantic proximity. A user querying "what are the security requirements for deploying a new application?" may receive thousands of matching documents, ranked by keyword frequency, with no mechanism for synthesizing a coherent, sourced answer across the relevant subset.

This thesis originates in a concrete instantiation of this problem. The \gls{dnsi} — the information systems directorate of a large organization — maintains an extensive documentation corpus on its SharePoint platform (SharePoint \gls{cerp}). This corpus includes technical architecture documents, security policies, operational procedures, and project deliverables. Business units within the DNSI face daily challenges in locating, interpreting, and synthesizing information from this corpus, a friction that imposes measurable operational and decision-making costs.

%------------------------------------------------------------------------------
\section{The Promise and Limits of Large Language Models}
\label{sec:intro-llm}
%------------------------------------------------------------------------------

The emergence of Large Language Models — neural architectures trained on massive corpora to generate coherent, contextually appropriate text \citep{Brown2020, Chowdhery2022} — has opened a fundamentally new interface for knowledge access. An LLM can, in principle, receive a natural language query and produce a synthesized, readable, and structurally coherent response that directly addresses the user's need, without requiring the user to formulate Boolean search expressions or manually read through ranked document lists.

The practical deployment of LLMs in enterprise environments, however, is blocked by three fundamental and interdependent limitations:

\paragraph{Hallucination.} LLMs are generative models trained to produce plausible continuations of text, not to retrieve verified facts. They will generate confident, fluent responses that are factually incorrect — a phenomenon termed hallucination \citep{Maynez2020, Ji2023}. In a consumer context, this is problematic. In an enterprise context, where responses may inform operational decisions, security configurations, or regulatory compliance positions, it is unacceptable.

\paragraph{Knowledge Cutoff and Private Data.} LLMs are trained on static snapshots of public data up to a fixed point in time. They have no knowledge of an organization's private documentation, internal processes, post-training events, or domain-specific terminology. A model trained on public internet data cannot answer questions about the DNSI's internal security policy, because that policy was never in its training corpus. This is not a solvable problem through better model training — it is structural.

\paragraph{Traceability and Access Control.} Enterprise environments require that responses be attributable to specific, authoritative sources, and that access to sensitive documentation be governed by the organizational permission model. A standalone LLM provides neither: it generates responses without citing sources, and it applies no access control — the same model, with the same weights, would generate the same response regardless of whether the querying user is authorized to access the underlying information.

These three limitations converge to make direct LLM deployment in enterprise knowledge management contexts not merely suboptimal but operationally and legally untenable in most regulated environments.

%------------------------------------------------------------------------------
\section{Retrieval-Augmented Generation as a Paradigm Shift}
\label{sec:intro-rag}
%------------------------------------------------------------------------------

The Retrieval-Augmented Generation paradigm, introduced by \citet{Lewis2020} at NeurIPS 2020, offers a compelling architectural response to these limitations. Rather than relying exclusively on the parametric knowledge encoded in model weights, \gls{rag} retrieves relevant documents from an external knowledge base at query time, injects them as context into the LLM's input, and uses the model to synthesize a response grounded in the retrieved evidence.

This architectural choice addresses all three limitations simultaneously: the model's response is constrained by the retrieved context (reducing hallucination), the knowledge base can be populated with private, up-to-date organizational documents (addressing the knowledge cutoff), and the retrieval step can be restricted to documents the querying user is authorized to access (enforcing access control). Importantly, the knowledge base can be updated without retraining the model — a critical operational advantage in dynamic organizational environments where documentation evolves continuously.

Since Lewis et al.'s foundational work, the RAG paradigm has matured rapidly. \citet{Gao2024} survey over 100 RAG variants across three evolutionary generations: naive RAG (direct retrieval and generation), advanced RAG (pre-retrieval query optimization and post-retrieval reranking), and modular RAG (flexible pipeline architectures with specialized components for specific use cases). Each generation addresses limitations of the previous through increasingly sophisticated retrieval, context management, and generation strategies.

Despite this maturity at the architectural level, a critical infrastructure challenge remains largely unsolved: the problem of the data connector. Enterprise knowledge is stored in heterogeneous, siloed systems — SharePoint, Confluence, SAP, internal databases, email archives — each with its own API, authentication model, document format, and permission structure. Integrating RAG pipelines with these sources has historically required bespoke, hard-to-maintain connectors that introduce significant engineering overhead and security risk. This integration problem is the immediate motivation for the Model Context Protocol.

%------------------------------------------------------------------------------
\section{The Model Context Protocol: Standardizing the AI-Data Interface}
\label{sec:intro-mcp}
%------------------------------------------------------------------------------

In November 2024, Anthropic published the specification for the Model Context Protocol (MCP) \citep{Anthropic2024MCP} — an open standard designed to formalize the interface between AI reasoning systems and external data sources. MCP defines a JSON-RPC-based protocol through which AI clients can discover, connect to, and query data providers (MCP servers) through a unified, standardized interface, regardless of the underlying data source's native API.

The architectural significance of MCP is the introduction of a stable abstraction layer between the AI layer and the data layer. Before MCP, connecting an AI assistant to SharePoint required writing SharePoint-specific code directly in the AI pipeline — code that couples the pipeline to SharePoint's API versioning, authentication changes, and structural evolution. With MCP, the SharePoint-specific logic is encapsulated in an MCP server, and the AI pipeline communicates through the stable MCP protocol. Changes in the SharePoint API require only updating the MCP server, not the AI pipeline. This separation of concerns is architecturally analogous to the role ODBC played in standardizing database connectivity in the 1990s.

Beyond maintainability, MCP's architecture introduces properties particularly relevant to enterprise environments: the server-side context management allows access control to be enforced at the data boundary (before context is passed to the AI model), making permission enforcement structural rather than advisory. The protocol's tool-calling semantics also enable richer interactions than simple document retrieval — an MCP server can expose structured operations (search, filter, navigate document hierarchies) that a naive retrieval connector cannot.

However, MCP is a recent standard. Its practical value in production enterprise deployments — particularly regarding the claimed benefits of security, interoperability, and maintainability — has not yet been rigorously evaluated in a published academic study. This empirical gap constitutes the primary scientific motivation for this thesis.

%------------------------------------------------------------------------------
\section{Research Problem and Questions}
\label{sec:intro-problem}
%------------------------------------------------------------------------------

The preceding analysis converges on the following primary research question:

\begin{quote}
\textbf{RQ1 (Primary):} \textit{To what extent does a Retrieval-Augmented Generation architecture grounded in the Model Context Protocol enable more reliable, secure, and traceable access to enterprise knowledge stored in SharePoint, compared to a standalone LLM and a naive RAG implementation?}
\end{quote}

This primary question is operationalized through four subsidiary research questions:

\begin{itemize}
  \item \textbf{RQ2 (Retrieval):} What retrieval configuration (chunking strategy, embedding model, similarity metric, re-ranking approach) maximizes retrieval precision and recall on a domain-specific SharePoint document corpus?

  \item \textbf{RQ3 (Protocol):} What specific security, interoperability, and maintainability properties does MCP deliver in an enterprise SharePoint integration context, and how do these compare to alternative integration approaches?

  \item \textbf{RQ4 (Performance):} What are the latency, throughput, and scalability characteristics of the RAG+MCP system under realistic enterprise workload conditions?

  \item \textbf{RQ5 (Acceptance):} How do end-users perceive the usefulness, trustworthiness, and ease of use of the AntiGravity assistant compared to their current SharePoint search workflow?
\end{itemize}

%------------------------------------------------------------------------------
\section{Research Approach and Contributions}
\label{sec:intro-contributions}
%------------------------------------------------------------------------------

This thesis adopts the \textbf{Design Science Research} (DSR) paradigm \citep{Hevner2004}, which frames research as the construction and rigorous evaluation of a purposeful artifact. The artifact — the AntiGravity RAG+MCP system — constitutes the primary research output, while the evaluation of that artifact against a real organizational problem constitutes the primary scientific contribution.

The thesis makes the following contributions:

\paragraph{Scientific Contribution 1 (SC1):} A formal comparative evaluation framework for RAG+MCP architectures in enterprise settings, operationalized through a controlled experiment on a real SharePoint corpus with three conditions: (A) standalone LLM, (B) naive RAG, (C) RAG+MCP. This is, to the best of our knowledge, the first published empirical study to evaluate MCP-grounded RAG against baseline architectures on a private enterprise document corpus.

\paragraph{Scientific Contribution 2 (SC2):} A reference architecture for secure, \gls{acl}-aware RAG deployment in enterprise environments, including design patterns for MCP-compatible SharePoint connectors and a formal analysis of the security properties of the proposed architecture.

\paragraph{Scientific Contribution 3 (SC3):} An application of Design Science Research to the evaluation of enterprise AI systems, demonstrating how DSR can structure both the artifact construction and the artifact evaluation in a single methodologically coherent framework.

\paragraph{Technical Contribution 1 (TC1):} An MCP-compatible SharePoint connector that extracts documents from SharePoint, respects ACLs, and exposes the corpus through the MCP interface.

\paragraph{Technical Contribution 2 (TC2):} A domain-specific benchmark dataset of question-answer pairs derived from the DNSI SharePoint corpus, suitable for evaluating any conversational assistant on enterprise business documents.

%------------------------------------------------------------------------------
\section{Scope and Boundaries}
\label{sec:intro-scope}
%------------------------------------------------------------------------------

This thesis addresses the design, implementation, and evaluation of a conversational assistant for enterprise document access. The following boundaries apply:

\textbf{In scope:} Design of the AntiGravity architecture; implementation of the RAG+MCP pipeline; evaluation against three experimental conditions; retrieval, generation, system, and user evaluation; security architecture analysis; DSR-grounded methodology.

\textbf{Out of scope:} Training or fine-tuning of language models; multi-tenant deployment; real-time SharePoint delta synchronization; extension to data sources beyond SharePoint; production-scale infrastructure optimization. These are identified as directions for future work in Chapter~\ref{chap:discussion}.

%------------------------------------------------------------------------------
\section{Thesis Structure}
\label{sec:intro-structure}
%------------------------------------------------------------------------------

The remainder of this thesis is organized as follows:

\textbf{Chapter~\ref{chap:background}} establishes the theoretical and empirical foundations of the work through a comprehensive review of the literature on LLMs, RAG architectures, vector search, enterprise knowledge management, the Model Context Protocol, and enterprise AI security. It concludes by identifying the research gap this thesis fills.

\textbf{Chapter~\ref{chap:methodology}} formalizes the research methodology, grounding the approach in Design Science Research, and specifies the complete evaluation protocol for each of the four evaluation dimensions.

\textbf{Chapter~\ref{chap:architecture}} presents the system design and architecture of the AntiGravity assistant, documenting the design decisions made at each architectural layer, the alternatives considered and rejected, and the trade-offs accepted.

\textbf{Chapter~\ref{chap:implementation}} describes the technical implementation of the system, documenting the engineering challenges encountered and the solutions developed.

\textbf{Chapter~\ref{chap:evaluation}} presents the experimental results across all four evaluation dimensions, tests the six research hypotheses, and provides detailed statistical analysis of the findings.

\textbf{Chapter~\ref{chap:discussion}} provides a critical discussion of the results, analyzing strengths, weaknesses, limitations, deployment challenges, security gaps, and the generalizability of the findings.

\textbf{Chapter~\ref{chap:conclusion}} synthesizes the contributions of the thesis, provides explicit answers to each research question, and outlines the directions for future work, including a proposed research agenda for CIFRE doctoral continuation.

%==============================================================================
% CHAPTER 2 — BACKGROUND AND STATE OF THE ART
%==============================================================================
\chapter{Background and State of the Art}
\label{chap:background}

This chapter establishes the theoretical and empirical foundations required to understand, position, and evaluate the contributions of this thesis. We survey the relevant literature across six interconnected domains: Large Language Models and their limitations (Section~\ref{sec:bg-llm}); Retrieval-Augmented Generation architectures (Section~\ref{sec:bg-rag}); vector search and semantic indexing (Section~\ref{sec:bg-vector}); enterprise knowledge management and SharePoint (Section~\ref{sec:bg-enterprise}); the Model Context Protocol (Section~\ref{sec:bg-mcp}); security in enterprise AI systems (Section~\ref{sec:bg-security}); and RAG evaluation frameworks (Section~\ref{sec:bg-evaluation}). Each section concludes with a critical assessment of the current state of the art and the limitations this thesis addresses. Section~\ref{sec:bg-gap} synthesizes the literature into an explicit research gap statement.

%------------------------------------------------------------------------------
\section{Large Language Models: Capabilities and Limitations}
\label{sec:bg-llm}
%------------------------------------------------------------------------------

\subsection{The Transformer Architecture}

The modern era of Large Language Models originates with the Transformer architecture introduced by \citet{Vaswani2017} in their seminal paper "Attention Is All You Need." The Transformer replaces the sequential computation of recurrent networks with a self-attention mechanism that computes, for each token in a sequence, a weighted representation of all other tokens simultaneously. This design enables parallelization during training — a property that, combined with scale, proved to be the foundation of modern language models.

The self-attention mechanism can be formalized as:
\begin{equation}
  \text{Attention}(Q, K, V) = \text{softmax}\left(\frac{QK^\top}{\sqrt{d_k}}\right)V
  \label{eq:attention}
\end{equation}
where $Q$, $K$, and $V$ are query, key, and value matrices respectively, and $d_k$ is the dimensionality of the key vectors. The scaling factor $\sqrt{d_k}$ prevents vanishing gradients in high-dimensional spaces.

Modern large-scale language models — including GPT-4 \citep{OpenAI2023}, Mistral 7B \citep{Jiang2023}, Llama 2/3 \citep{Touvron2023}, and Claude \citep{Anthropic2023} — are decoder-only or encoder-decoder Transformers trained via next-token prediction on datasets of trillions of tokens. The emergent capabilities observed at scale — instruction following, chain-of-thought reasoning, code generation, multi-step problem solving — are not explicitly programmed but arise from the optimization process applied at sufficient scale \citep{Wei2022}.

\subsection{Parametric vs Non-Parametric Knowledge}

A foundational distinction for understanding LLM limitations is between \textit{parametric} and \textit{non-parametric} knowledge \citep{Lewis2020}. Parametric knowledge is information encoded in a model's weights during training — it is fixed, distributed, implicit, and inaccessible without inference. Non-parametric knowledge is information stored in explicit, updatable external databases — structured, addressable, and modifiable without retraining.

LLMs are parametic knowledge stores. Everything they "know" is compressed into billions of floating-point weights through the training process. This has several important implications:

\begin{itemize}
  \item \textbf{Knowledge is probabilistic, not factual.} The model learns statistical patterns in text, not ground truth facts. Information that is rare, contradictory, or absent in the training corpus will be poorly represented or incorrectly recalled.
  \item \textbf{Knowledge is static.} Once training is complete, the model's knowledge is frozen. Events after the training cutoff, or information not present in training data, are genuinely unknown to the model.
  \item \textbf{Knowledge is opaque.} There is no mechanism for tracing a model's assertion to a specific source document. The model cannot say "I know this because I read it in document X" — it can only generate what its weights make probable.
\end{itemize}

\subsection{The Hallucination Problem}

Hallucination — the generation of confident, fluent, but factually incorrect text — is the most operationally significant limitation of LLMs for enterprise deployment. \citet{Ji2023} provide a taxonomy of hallucination types, distinguishing between intrinsic hallucinations (contradictions of source material) and extrinsic hallucinations (unverifiable claims absent from any source). \citet{Maynez2020} demonstrate that even state-of-the-art summarization models hallucinate in 30\% or more of generated summaries.

The underlying mechanism of hallucination is partially understood. Models optimize for textual fluency and statistical plausibility, not factual accuracy. When asked about information that is rare, ambiguous, or absent in training data, the model continues to generate plausible text — because that is what it was trained to do — regardless of whether the generated content is grounded in fact \citep{Rawte2023}.

For enterprise applications, hallucination is not merely a quality problem — it is a trust and liability problem. A conversational assistant that occasionally fabricates security policy requirements, invents technical specifications, or misrepresents regulatory obligations can cause real operational harm. Any architecture deployed in an enterprise knowledge management context must provide structural mechanisms for constraining model generation to verified, retrievable source material.

\subsection{Critical Assessment}

Despite their remarkable capabilities, standalone LLMs are insufficient for enterprise knowledge access for three reasons that cannot be resolved through better model architecture or larger training data: (1) private organizational knowledge will never appear in public training corpora; (2) static training weights cannot reflect dynamic organizational documentation; and (3) parametric generation provides no mechanism for access control or source attribution. The RAG architecture addresses all three limitations. We now examine it in detail.

%------------------------------------------------------------------------------
\section{Retrieval-Augmented Generation}
\label{sec:bg-rag}
%------------------------------------------------------------------------------

\subsection{Foundational Architecture}

Retrieval-Augmented Generation was formalized by \citet{Lewis2020} as a general-purpose architecture for knowledge-intensive NLP tasks. The original formulation combines a parametric generative model (such as BART or GPT) with a non-parametric dense retrieval system (DPR — Dense Passage Retrieval), end-to-end trained to improve performance on open-domain question answering benchmarks.

The RAG inference process follows a three-stage pipeline:

\begin{enumerate}
  \item \textbf{Query encoding:} The user's query $q$ is encoded by a query encoder $E_q$ into a dense vector representation $\mathbf{q} = E_q(q) \in \mathbb{R}^d$.

  \item \textbf{Retrieval:} The query vector $\mathbf{q}$ is used to retrieve the $k$ most semantically similar document chunks from an indexed corpus. Formally:
  \begin{equation}
    D_k = \arg\max_{d \in \mathcal{D}} \text{sim}(\mathbf{q}, E_d(d))
    \label{eq:retrieval}
  \end{equation}
  where $\mathcal{D}$ is the document corpus, $E_d$ is the document encoder, and $\text{sim}(\cdot, \cdot)$ is a similarity function (typically dot product or cosine similarity).

  \item \textbf{Generation:} The retrieved documents $D_k$ are concatenated with the original query and passed to the generative model $G$ to produce the final response:
  \begin{equation}
    y = G(q, D_k)
    \label{eq:generation}
  \end{equation}
\end{enumerate}

The central insight of this architecture is that the retrieval step \textit{grounds} the generation step: the model's output is constrained by, and should be attributable to, the retrieved context. This grounding mechanism directly mitigates hallucination by providing the model with relevant, specific, authoritative information at generation time.

\subsection{RAG Variants and Evolution}

The three-year period following Lewis et al.'s foundational work has seen rapid architectural evolution. \citet{Gao2024} organize this evolution into three generations:

\paragraph{Naive RAG.} The direct implementation of the Lewis et al. pipeline: retrieve, then generate. Simple and efficient, but susceptible to retrieval failure (the generation can only be as good as what was retrieved) and context stuffing (injecting low-quality retrieved passages that confuse the model rather than grounding it).

\paragraph{Advanced RAG.} Introduces improvements at both the pre-retrieval and post-retrieval stages. Pre-retrieval techniques include query rewriting \citep{Ma2023}, HyDE (Hypothetical Document Embeddings) \citep{Gao2022HyDE}, and step-back prompting. Post-retrieval techniques include re-ranking (filtering and reordering retrieved passages by relevance to the query), context compression (removing low-information content from retrieved passages), and fusion retrieval (combining dense and sparse retrieval signals).

\paragraph{Modular RAG.} Generalizes the pipeline into a configurable sequence of specialized modules — query classification, routing, retrieval, re-ranking, fusion, generation, and evaluation — that can be independently optimized and combined. Frameworks such as LangChain \citep{Chase2022} and LlamaIndex \citep{Liu2022} provide infrastructure for constructing modular RAG pipelines.

\subsection{Limitations of Current RAG Implementations}

Despite its maturity, current RAG research and tooling presents three significant limitations for enterprise deployment:

\textbf{Connector brittleness.} Most RAG implementations assume a static, pre-processed document corpus. Enterprise knowledge is dynamic: documents are added, modified, and deleted continuously. Maintaining a synchronized index requires reliable, efficient data connectors, which existing frameworks handle inconsistently.

\textbf{Access control is an afterthought.} The dominant RAG frameworks (LangChain, LlamaIndex) do not provide native mechanisms for enforcing document-level access control. Implementing ACL-aware retrieval requires custom engineering, and the resulting solutions are typically tightly coupled to the specific data source's permission model.

\textbf{Connector standardization is absent.} Each data source (SharePoint, Confluence, S3, databases) requires a different connector, written against a different API, with different authentication semantics. This creates an O(n × m) integration problem that grows with the number of data sources and AI pipeline variants.

The Model Context Protocol directly targets this third limitation, and its architecture has implications for the first two as well.

%------------------------------------------------------------------------------
\section{Vector Search and Semantic Indexing}
\label{sec:bg-vector}
%------------------------------------------------------------------------------

\subsection{Dense Embeddings}

The retrieval step in a RAG pipeline depends on the ability to represent both queries and documents as dense vectors in a shared semantic space, such that similar meanings produce geometrically proximate representations. This requires an \textit{embedding model} — a neural encoder that maps variable-length text to fixed-dimensional vectors.

The field has converged on bi-encoder architectures trained with contrastive objectives \citep{Karpukhin2020DPR}. Models such as \texttt{sentence-transformers/all-MiniLM-L6-v2} \citep{Reimers2019}, \texttt{text-embedding-ada-002} (OpenAI), and \texttt{e5-large} \citep{Wang2022} produce embeddings that capture semantic similarity significantly better than lexical representations (TF-IDF, BM25) on most open-domain benchmarks. The trade-off between embedding quality and computational cost (embedding dimensionality, model size, inference speed) is a critical engineering decision in any RAG deployment.

\subsection{FAISS: Scalable Similarity Search}

The \textit{Facebook AI Similarity Search} (\gls{faiss}) library \citep{Johnson2019} provides state-of-the-art infrastructure for billion-scale approximate nearest neighbor (\gls{ann}) search in dense vector spaces. FAISS implements multiple index types with different precision-speed trade-offs:

\begin{table}[H]
\centering
\caption{FAISS index types and their trade-offs}
\label{tab:faiss-indexes}
\begin{tabular}{@{}llll@{}}
\toprule
\textbf{Index Type} & \textbf{Precision} & \textbf{Speed} & \textbf{Memory} \\
\midrule
\texttt{IndexFlatL2} & Exact (100\%) & Slow (O(n)) & High \\
\texttt{IndexIVFFlat} & High (95–99\%) & Fast & Moderate \\
\texttt{IndexHNSW} & High (95–99\%) & Very fast & High \\
\texttt{IndexIVFPQ} & Moderate (85–95\%) & Very fast & Low \\
\bottomrule
\end{tabular}
\end{table}

For the document scales typical of an enterprise SharePoint corpus (thousands to hundreds of thousands of documents), \texttt{IndexFlatL2} (exact search) is computationally feasible and maximally precise, while \texttt{IndexHNSW} provides the best speed-precision trade-off at larger scales. The optimal choice depends on corpus size, query latency requirements, and available memory — all of which must be measured empirically for the specific deployment context.

\subsection{Re-ranking}

Bi-encoder retrieval trades recall for precision: the shared embedding space, trained for broad semantic similarity, may retrieve documents that are topically related but not directly responsive to the specific query. Re-ranking addresses this by applying a more computationally expensive but more discriminative model to the top-$k$ retrieved candidates, reordering them by relevance to the specific query.

Cross-encoder architectures \citep{Nogueira2019}, which process the (query, document) pair jointly rather than independently, consistently outperform bi-encoder retrieval on precision metrics at the cost of higher inference latency. ColBERT \citep{Khattab2020} provides an intermediate approach using late interaction, achieving cross-encoder precision at closer to bi-encoder speed. The choice between these approaches is another deployment-critical engineering decision.

%------------------------------------------------------------------------------
\section{Enterprise Knowledge Management and SharePoint}
\label{sec:bg-enterprise}
%------------------------------------------------------------------------------

\subsection{Enterprise Knowledge Management}

Knowledge Management (KM) as a discipline addresses the creation, capture, storage, sharing, and application of organizational knowledge \citep{Alavi2001}. The field distinguishes between explicit knowledge (codified in documents, procedures, and systems) and tacit knowledge (embedded in individual expertise and practice). Enterprise AI systems such as the one developed in this thesis target explicit knowledge — the documented corpus accessible on platforms like SharePoint.

The organizational context matters for system design. Enterprises operate under regulatory constraints (data protection, information security, sector-specific compliance), organizational hierarchies (which map to document access permissions), and professional vocabulary (domain-specific terminology that generic embedding models may not represent well). A system designed for Wikipedia-scale open-domain question answering will not necessarily perform well on an enterprise corpus — a point that motivates the domain-specific evaluation protocol in this thesis.

\subsection{SharePoint Architecture and Security Model}

Microsoft SharePoint Online is built on a hierarchical content structure: Tenant → Site Collection → Site → Library → Folder → Document. Each level of this hierarchy can carry permission settings, which are either inherited (the default) or explicitly overridden. The SharePoint permission model comprises permission levels (Full Control, Design, Edit, Contribute, Read, View Only) assigned to security principals (users, Microsoft Entra ID groups, SharePoint groups) \citep{Microsoft2024}.

At the API level, SharePoint exposes two integration surfaces relevant to this thesis:

\begin{itemize}
  \item \textbf{Microsoft Graph API} \citep{MicrosoftGraph2024}: a REST API providing unified access to Microsoft 365 services, including SharePoint document libraries. Authentication uses OAuth 2.0 with delegated or application permissions.
  \item \textbf{SharePoint REST API}: the legacy SharePoint-specific API, still widely used for site-level operations.
\end{itemize}

The key security challenge for RAG-over-SharePoint is \textit{permission-aware retrieval}: the semantic search that identifies relevant document chunks must respect the access rights of the querying user. A naive implementation — building a FAISS index over all documents, regardless of permissions, and then checking permissions at display time — would expose the existence of unauthorized documents to the user (through responses that implicitly reference them) and potentially leak information if the authorization check fails or is bypassed. A secure implementation must enforce permissions at the \textit{retrieval} layer, not the display layer. Designing this enforcement mechanism is a primary technical contribution of this thesis.

%------------------------------------------------------------------------------
\section{The Model Context Protocol}
\label{sec:bg-mcp}
%------------------------------------------------------------------------------

\subsection{Protocol Overview}

The Model Context Protocol (MCP) \citep{Anthropic2024MCP} was publicly released by Anthropic in November 2024 as an open standard for connecting AI systems to external data and tool providers. As of mid-2026, MCP has been adopted by major AI platforms including Claude (Anthropic), Copilot (Microsoft), and numerous enterprise AI frameworks, positioning it as a de facto standard for AI-data integration.

MCP defines a client-server architecture with three roles:

\begin{itemize}
  \item \textbf{MCP Host:} The AI application or agent (in our case, the AntiGravity conversational client) that initiates connections and uses capabilities.
  \item \textbf{MCP Server:} A lightweight process that exposes data and tools to the host through the MCP protocol (in our case, the SharePoint context server).
  \item \textbf{MCP Transport:} The communication layer between host and server, supporting stdio (for local processes) and HTTP with Server-Sent Events (for networked deployments).
\end{itemize}

MCP defines three categories of server-exposed capabilities:

\begin{itemize}
  \item \textbf{Resources:} Data that the AI can read — documents, files, database records, API responses.
  \item \textbf{Tools:} Functions the AI can call — search, filter, write, API operations.
  \item \textbf{Prompts:} Reusable prompt templates that the server exposes to the AI.
\end{itemize}

\subsection{MCP vs Alternative Integration Approaches}

The value of MCP can only be properly assessed in comparison to the alternatives it is designed to replace. Table~\ref{tab:mcp-comparison} provides a structured comparison.

\begin{table}[H]
\centering
\caption{Comparison of AI-data integration approaches}
\label{tab:mcp-comparison}
\begin{tabularx}{\textwidth}{@{}lXXX@{}}
\toprule
\textbf{Dimension} & \textbf{Direct API} & \textbf{LangChain / LlamaIndex} & \textbf{MCP} \\
\midrule
Standardization & None (source-specific) & Framework-specific & Open protocol standard \\
Connector reuse & None & Framework ecosystem & Cross-framework (any MCP client) \\
Security boundary & In-pipeline code & In-pipeline code & At server boundary (structural) \\
Access control & Ad hoc & Ad hoc & Enforced at server level \\
Maintainability & Low (coupled) & Medium (framework dep.) & High (interface stability) \\
Overhead & Minimal & Low & Low-to-moderate (protocol) \\
Maturity (2026) & High & High & Emerging but fast-growing \\
Enterprise adoption & High & Medium & Growing (Microsoft, Anthropic) \\
\bottomrule
\end{tabularx}
\end{table}

The comparison reveals MCP's primary advantage: structural security enforcement and cross-framework connector reuse. Its primary disadvantage — lower maturity and higher protocol overhead compared to direct API integration — is the key empirical question this thesis addresses.

\subsection{MCP in Enterprise Deployments: Open Questions}

Despite MCP's architectural promise, several questions remain unanswered in the literature:

\begin{enumerate}
  \item What latency overhead does the MCP protocol layer introduce in practice, and is this overhead acceptable in enterprise SLA contexts?
  \item How does MCP's structured context management compare to direct context injection in terms of retrieval quality and generation faithfulness?
  \item Can MCP's server-side security enforcement reliably enforce SharePoint ACLs without performance degradation?
  \item What is the operational complexity of deploying an MCP server in a corporate network environment with proxy constraints?
\end{enumerate}

These are the questions that Chapter~\ref{chap:evaluation} empirically addresses.

%------------------------------------------------------------------------------
\section{Security in Enterprise AI Systems}
\label{sec:bg-security}
%------------------------------------------------------------------------------

\subsection{Information Security Requirements}

Enterprise AI systems operate within information security frameworks that impose binding requirements on data handling, access control, and audit logging. The ISO/IEC 27001 standard \citep{ISO27001} defines a comprehensive Information Security Management System (ISMS) framework whose principles are directly applicable to AI system design: access control (Annex A.9), cryptography (A.10), operations security (A.12), and supplier relationships (A.15) are all implicated in the deployment of a conversational AI system over enterprise document corpora.

The specific security concerns for a RAG-over-SharePoint system include:

\paragraph{Document access leakage.} If the retrieval system retrieves and injects into the LLM prompt a document the querying user is not authorized to access, the LLM may reflect information from that document in its response — even if the document is not explicitly cited. This indirect leakage is subtle and potentially more dangerous than direct access because it may not be detected by conventional access control audit tools.

\paragraph{Prompt injection.} An adversarial actor with write access to SharePoint could plant documents containing prompt injection instructions — text designed to manipulate the LLM's behavior when retrieved as context \citep{Perez2022}. In an enterprise setting, this attack vector requires a threat model that accounts for internal adversaries.

\paragraph{Model output as an attack surface.} LLM-generated responses that are displayed to users create a novel attack surface. If the model can be manipulated to generate links, JavaScript, or other active content, the UI layer must sanitize model outputs.

\subsection{Authentication: OAuth 2.0 and SSO}

Secure access to SharePoint via the Microsoft Graph API requires OAuth 2.0 authentication, with either delegated permissions (the application acts on behalf of a specific user, inheriting their permissions) or application permissions (the application has its own permissions, granted by an administrator). For an ACL-aware RAG system, \textit{delegated permissions are mandatory} — only delegated permissions allow the system to enforce the specific document access rights of the querying user, rather than the broader application-level permissions.

In practice, this means that each user query must be executed under the user's OAuth token — a requirement that introduces session management complexity and latency overhead, but is non-negotiable from a security correctness standpoint.

%------------------------------------------------------------------------------
\section{RAG Evaluation Frameworks}
\label{sec:bg-evaluation}
%------------------------------------------------------------------------------

\subsection{The RAGAS Framework}

The evaluation of RAG systems requires metrics that assess both retrieval quality and generation quality, independently and jointly. The RAGAS framework \citep{Es2023RAGAS} provides an automated evaluation suite specifically designed for RAG pipelines, computing four primary metrics:

\begin{itemize}
  \item \textbf{Faithfulness:} The proportion of claims in the generated response that are grounded in the retrieved context. A faithful response makes only claims that can be directly inferred from the provided context. Computed as:
  \begin{equation}
    \text{Faithfulness} = \frac{|\text{claims in response supported by context}|}{|\text{total claims in response}|}
    \label{eq:faithfulness}
  \end{equation}

  \item \textbf{Answer Relevancy:} The semantic similarity between the generated response and the original question, measuring whether the response addresses what was asked.

  \item \textbf{Context Precision:} The proportion of the retrieved context that is actually relevant to the question, measuring retrieval efficiency.

  \item \textbf{Context Recall:} The proportion of relevant information (as judged against a ground-truth answer) that is present in the retrieved context.
\end{itemize}

\subsection{Retrieval Evaluation Metrics}

Classical information retrieval metrics provide complementary evaluation of the retrieval stage:

\begin{itemize}
  \item \textbf{Precision@K:} The proportion of the top-$K$ retrieved documents that are relevant. Precision@5 is the most commonly reported metric in RAG evaluation studies.
  \item \textbf{Recall@K:} The proportion of all relevant documents that appear in the top-$K$ retrieved results.
  \item \textbf{Mean Reciprocal Rank (MRR):} The average of the reciprocal ranks of the first relevant document across queries. Captures the position of the first correct result.
  \item \textbf{Normalized Discounted Cumulative Gain (nDCG@K):} A rank-weighted precision metric that accounts for the position of relevant documents within the retrieved list.
\end{itemize}

\subsection{User Evaluation: SUS and TAM}

Technical retrieval and generation metrics are necessary but insufficient for evaluating an enterprise system. The system must also be usable and accepted by its target users. Two validated instruments are employed in this thesis:

The \textbf{System Usability Scale} (SUS) \citep{Brooke1996SUS} is a 10-item Likert questionnaire that produces a composite usability score from 0 to 100. SUS scores above 68 are considered above average, above 80 are considered excellent. SUS is technology-agnostic, widely validated, and appropriate for comparative usability assessment.

The \textbf{Technology Acceptance Model} (TAM) \citep{Davis1989TAM} models user adoption through two primary constructs: perceived usefulness (PU) and perceived ease of use (PEOU). TAM has been extensively validated in enterprise software deployment contexts and provides theoretical grounding for interpreting user acceptance of AI-powered tools.

%------------------------------------------------------------------------------
\section{Research Gap and Thesis Positioning}
\label{sec:bg-gap}
%------------------------------------------------------------------------------

The preceding review of the literature identifies a specific and well-defined research gap. Table~\ref{tab:research-gap} summarizes this positioning.

\begin{table}[H]
\centering
\caption{Research gap: what existing work covers vs what this thesis contributes}
\label{tab:research-gap}
\begin{tabularx}{\textwidth}{@{}lXX@{}}
\toprule
\textbf{Dimension} & \textbf{Existing Literature} & \textbf{This Thesis} \\
\midrule
Knowledge base & Wikipedia, open web, academic QA & Private enterprise SharePoint corpus \\
Integration protocol & Ad-hoc connectors / LangChain & Model Context Protocol (standardized) \\
Access control & Generally absent or post-hoc & ACL-enforced at retrieval layer (structural) \\
Comparative evaluation & LLM vs RAG (general) & LLM vs RAG vs RAG+MCP (enterprise) \\
Evaluation scope & Retrieval or generation & Retrieval + generation + system + user \\
Industrial context & Academic benchmarks & Real DNSI production constraints \\
DSR methodology & Rarely applied to RAG systems & Full DSR cycle applied \\
\bottomrule
\end{tabularx}
\end{table}

The specific combination of (1) MCP as the integration protocol, (2) ACL-aware retrieval enforcement, (3) private enterprise corpus, and (4) comprehensive four-dimensional evaluation has not appeared in any prior published work as of the writing of this thesis. This combination constitutes the primary scientific novelty of the thesis and the direct motivation for the research questions formulated in Chapter~\ref{chap:introduction}.

\bigskip

\noindent\textbf{Chapter summary.} This chapter has established that: (1) standalone LLMs are structurally insufficient for enterprise knowledge access due to hallucination, knowledge cutoff, and the absence of access control; (2) RAG architectures address these limitations but existing implementations fail to provide standardized, secure integration with enterprise data sources; (3) FAISS-based dense retrieval with re-ranking represents the state of the art for scalable semantic search; (4) the Model Context Protocol provides a promising but empirically unevaluated standardization layer for AI-enterprise data integration; and (5) comprehensive RAG evaluation requires four complementary dimensions — retrieval, generation, system, and user. The following chapter describes the methodology adopted to address these gaps.

%==============================================================================
% CHAPTER 3 — RESEARCH METHODOLOGY
%==============================================================================
\chapter{Research Methodology}
\label{chap:methodology}

%------------------------------------------------------------------------------
\section{Choice of Research Paradigm}
\label{sec:method-paradigm}
%------------------------------------------------------------------------------

The research question at the center of this mémoire — how does a RAG+MCP architecture compare to alternatives for enterprise knowledge access? — is simultaneously a design question (how should the system be built?) and an evaluation question (how does the system perform?). This dual nature requires a research paradigm that can accommodate both artifact construction and empirical evaluation within a single coherent framework.

We adopt \textbf{Design Science Research} (DSR), as formalized by \citet{Hevner2004} and further developed by \citet{Peffers2007}. DSR is a research paradigm in Information Systems and Engineering that frames research as the creation of purposeful \textit{artifacts} — constructs, models, methods, and instantiations — that solve real-world problems. The artifact is not merely a deliverable; it is the primary vehicle through which knowledge is generated and communicated. DSR is appropriate for this mémoire for three reasons:

\begin{enumerate}
  \item The primary output is a designed system artifact (AntiGravity) whose design decisions constitute the research contribution.
  \item The artifact is evaluated against a real organizational problem in its natural context (DNSI / SharePoint CERP), not against a synthetic benchmark.
  \item The research is expected to produce both prescriptive knowledge (design principles that generalize beyond the specific artifact) and descriptive knowledge (empirical evaluation results that inform future practitioners).
\end{enumerate}

Alternative paradigms were considered and rejected. A pure experiment approach would have treated the artifact as a black box and focused exclusively on outcome metrics, neglecting the architectural and design contributions. A case study approach would have produced rich qualitative understanding but insufficient empirical rigor for the quantitative evaluation required by the research questions. Action research would have entailed a more iterative, participatory process with DNSI stakeholders than was operationally feasible.

%------------------------------------------------------------------------------
\section{DSR Cycles Applied to This PFE}
\label{sec:method-cycles}
%------------------------------------------------------------------------------

\citet{Hevner2007} describe DSR as operating through three interconnected cycles. Table~\ref{tab:dsr-cycles} maps each cycle to the specific activities of this mémoire.

\begin{table}[H]
\centering
\caption{Design Science Research cycles applied to this PFE}
\label{tab:dsr-cycles}
\begin{tabularx}{\textwidth}{@{}lXX@{}}
\toprule
\textbf{DSR Cycle} & \textbf{General Definition} & \textbf{Application in This PFE} \\
\midrule
Relevance Cycle & Connects the research to the real-world context that motivates the problem & DNSI information access problem; SharePoint as knowledge repository; operational friction analysis \\
Design Cycle & Iterative construction and evaluation of the artifact & AntiGravity architecture design; chunking strategy selection; retrieval optimization; UI design \\
Rigor Cycle & Connects research to the knowledge base of existing theory and method & RAG literature (Lewis 2020, Gao 2024); MCP specification; DSR methodology; evaluation frameworks (RAGAS, SUS, TAM) \\
\bottomrule
\end{tabularx}
\end{table}

%------------------------------------------------------------------------------
\section{Research Design Overview}
\label{sec:method-design}
%------------------------------------------------------------------------------

The research proceeds through five phases, as illustrated in Figure~\ref{fig:research-design}. Each phase has clearly defined inputs, outputs, and validation criteria.

\begin{figure}[H]
\centering
\begin{tikzpicture}[
  phasebox/.style={draw, rounded corners=2pt, text width=0.84\textwidth,
                   align=center, inner ysep=6pt, font=\small, fill=lightgray},
  flow/.style={-{Stealth[length=2.5mm]}, thick},
  node distance=0.4cm
]
\node[phasebox] (p1) {\textbf{Phase 1 --- Problem Formulation}\\[1pt]
  {\footnotesize\itshape DNSI organizational context, literature review
   $\;\Rightarrow\;$ research questions, hypotheses, evaluation protocol}};
\node[phasebox, below=of p1] (p2) {\textbf{Phase 2 --- Artifact Design}\\[1pt]
  {\footnotesize\itshape Requirements analysis, literature on RAG/MCP/security
   $\;\Rightarrow\;$ system architecture, design decisions}};
\node[phasebox, below=of p2] (p3) {\textbf{Phase 3 --- Artifact Construction}\\[1pt]
  {\footnotesize\itshape Architecture, technology stack
   $\;\Rightarrow\;$ working AntiGravity implementation}};
\node[phasebox, below=of p3] (p4) {\textbf{Phase 4 --- Evaluation}\\[1pt]
  {\footnotesize\itshape Benchmark dataset, experimental protocol, user study
   $\;\Rightarrow\;$ quantitative and qualitative results}};
\node[phasebox, below=of p4] (p5) {\textbf{Phase 5 --- Analysis and Contribution}\\[1pt]
  {\footnotesize\itshape Evaluation results
   $\;\Rightarrow\;$ hypothesis testing, design principles, limitations, future work}};

\foreach \a/\b in {p1/p2, p2/p3, p3/p4, p4/p5}
  \draw[flow] (\a) -- (\b);

\draw[flow, dashed] (p4.east) -- ++(0.7,0) |- (p2.east)
  node[pos=0.25, right=2pt, font=\scriptsize\itshape, align=left, rotate=90, anchor=south]
  {design cycle iterations};
\end{tikzpicture}
\caption{Research design: five phases of the DSR process. Each phase transforms its inputs (left of the arrow) into validated outputs (right); the dashed edge marks the iterative design--evaluation loop of the DSR design cycle.}
\label{fig:research-design}
\end{figure}

%------------------------------------------------------------------------------
\section{Evaluation Methodology}
\label{sec:method-evaluation}
%------------------------------------------------------------------------------

The evaluation of the AntiGravity system follows a four-dimensional protocol designed to provide a comprehensive assessment of both the technical and human-centered aspects of the system. Each dimension is described below.

\subsection{Experimental Conditions}

The comparative evaluation uses three experimental conditions to answer RQ1:

\begin{itemize}
  \item \textbf{Condition A — LLM Alone:} Mistral 7B model accessed directly, without retrieval augmentation. Queries are answered from parametric knowledge only.
  \item \textbf{Condition B — Naive RAG:} The same Mistral 7B model with direct FAISS retrieval (no MCP layer, no re-ranking, standard top-$k$ injection).
  \item \textbf{Condition C — RAG+MCP (AntiGravity):} The full system as designed — MCP-mediated retrieval with ACL enforcement, re-ranking, and generation via Mistral 7B.
\end{itemize}

Comparing C to A isolates the total contribution of the RAG+MCP architecture. Comparing C to B isolates the specific contribution of the MCP layer and the re-ranking module. All conditions use the same underlying language model (Mistral 7B Instruct) with identical generation parameters (temperature=0.1, top\_p=0.9, max\_tokens=512, repetition\_penalty=1.1) to ensure that observed differences are attributable to the architecture, not to model variation.

\subsection{Benchmark Dataset Construction}

A domain-specific benchmark dataset is constructed from the DNSI SharePoint corpus following this protocol:

\begin{enumerate}
  \item \textbf{Document sampling:} A stratified sample of documents is drawn from the SharePoint corpus, covering the major document categories present (technical documentation, procedural guides, security policies, project reports).
  \item \textbf{Question generation:} For each sampled document, four types of questions are generated: (i) factoid questions with a single-document answer; (ii) multi-hop questions requiring synthesis across two documents; (iii) procedural questions requiring step-ordered instructions; (iv) unanswerable questions whose answer is not in the corpus (testing the graceful-refusal requirement FR08).
  \item \textbf{Ground-truth annotation:} Each question is annotated with (a) the correct answer, (b) the set of relevant document chunks (passage-level ground truth), and (c) an answerability label.
  \item \textbf{Quality control:} Two independent annotators evaluate each question-answer pair. Overall inter-annotator agreement is measured with Cohen's $\kappa$; individual pairs on which the annotators disagree are reviewed and revised, and an overall $\kappa < 0.6$ triggers a full re-annotation pass.
\end{enumerate}

The target benchmark size is $N = 50$ question-answer pairs (20 factoid, 15 multi-hop, 10 procedural, 5 unanswerable), which permits reporting metric distributions with bootstrapped 95\% confidence intervals while acknowledging the limited statistical power inherent to this proof-of-concept evaluation (Section~\ref{sec:method-validity}).

\subsection{Retrieval Evaluation Protocol}

Retrieval quality is measured against the ground-truth relevant chunks for each benchmark question. The following metrics are computed for each retrieval-based experimental condition (B and C) and for a BM25 lexical baseline included to situate dense retrieval against classical keyword search (H3):

\begin{align}
\text{Precision@K} &= \frac{|\text{Relevant} \cap \text{Retrieved}@K|}{K} \label{eq:precision}\\[6pt]
\text{Recall@K} &= \frac{|\text{Relevant} \cap \text{Retrieved}@K|}{|\text{Relevant}|} \label{eq:recall}\\[6pt]
\text{MRR} &= \frac{1}{|Q|} \sum_{i=1}^{|Q|} \frac{1}{\text{rank}_i} \label{eq:mrr}\\[6pt]
\text{nDCG@K} &= \frac{\text{DCG@K}}{\text{IDCG@K}}, \quad \text{where DCG@K} = \sum_{i=1}^{K} \frac{\text{rel}_i}{\log_2(i+1)} \label{eq:ndcg}
\end{align}

Metrics are reported for $K \in \{1, 3, 5, 10\}$. Statistical significance is assessed with a paired Wilcoxon signed-rank test ($\alpha = 0.05$) comparing Condition B (Naive RAG) to Condition C (RAG+MCP).

\subsection{Generation Evaluation Protocol}

Generation quality is assessed across all three conditions using a combination of automated metrics (RAGAS) and human evaluation:

\paragraph{Automated metrics (RAGAS):}
\begin{itemize}
  \item Faithfulness: proportion of response claims grounded in retrieved context (Eq.~\ref{eq:faithfulness})
  \item Answer Relevancy: semantic similarity between response and ground-truth answer
  \item Context Precision and Context Recall (conditions B and C only)
\end{itemize}

\paragraph{Hallucination rate (human annotation):}
Each generated response is manually evaluated by two independent annotators for the presence of hallucinated claims (claims that contradict or are absent from the source documents). The hallucination rate is defined as:
\begin{equation}
  \text{Hallucination Rate} = \frac{|\text{Responses with} \geq 1 \text{ hallucinated claim}|}{|Q|}
  \label{eq:hallucination}
\end{equation}

All generation experiments are run with the system's production decoding configuration (temperature=0.1, fixed random seed) to ensure reproducibility. The difference in hallucination rate between paired conditions is assessed with McNemar's test on the per-query annotations. To assess variance introduced by prompt formulation, each query is run with two prompt variants and the metric range is reported alongside the mean.

\subsection{System Performance Evaluation Protocol}

The system is instrumented with timing middleware at the following pipeline boundaries: (1) query receipt; (2) query encoding completion; (3) FAISS retrieval completion; (4) re-ranking completion; (5) MCP context assembly completion; (6) LLM generation completion (local Mistral 7B); (7) response delivery.

Performance metrics are collected over $N=200$ representative queries executed in sequence:

\begin{itemize}
  \item End-to-end latency: P50, P95, P99 (milliseconds)
  \item Component-level latency breakdown: encoding, retrieval, re-ranking, MCP overhead, generation
  \item Throughput: queries per second under simulated concurrent load (5, 10, 20 concurrent users)
  \item Memory footprint: FAISS index resident memory, model memory
\end{itemize}

The MCP protocol overhead is isolated by comparing equivalent retrieval operations executed via the MCP server versus direct API calls.

\subsection{User Evaluation Protocol}

The user study is conducted with DNSI employees (target $N \geq 10$) following a structured protocol:

\begin{enumerate}
  \item \textbf{Task-based evaluation:} Each participant completes five representative business tasks using (a) the current SharePoint search interface and (b) the AntiGravity assistant. Task order is counterbalanced. Performance is measured by task completion rate and time-on-task.
  \item \textbf{SUS questionnaire:} After each system, participants complete the 10-item SUS questionnaire.
  \item \textbf{TAM questionnaire:} Participants complete a validated TAM instrument measuring perceived usefulness and perceived ease of use, administered for both systems so that AntiGravity can be compared to the incumbent SharePoint search interface.
  \item \textbf{Semi-structured interview:} A 15-minute qualitative interview explores trust, limitations perceived, and adoption barriers.
\end{enumerate}

Given the constraint of a small sample ($N < 30$), SUS and TAM scores are reported with confidence intervals using bootstrapping. Qualitative interview data is analyzed thematically \citep{Braun2006}.

%------------------------------------------------------------------------------
\subsection{Research Hypotheses}
\label{sec:method-hypotheses}

The following six hypotheses operationalize the research questions. Each hypothesis is evaluated in Chapter~\ref{chap:evaluation} against the corresponding metric.

\begin{itemize}
\item \textbf{H1:} The RAG+MCP architecture (Condition C) produces a statistically significantly lower hallucination rate than the standalone LLM (Condition A) on the DNSI benchmark.
\item \textbf{H2:} The RAG+MCP architecture achieves a RAGAS Faithfulness score $\geq 0.85$ on the benchmark (NFR10).
\item \textbf{H3:} Dense semantic retrieval (Condition B) achieves higher Precision@5 than BM25 lexical search on the DNSI corpus.
\item \textbf{H4:} The MCP protocol overhead does not exceed 200~ms at P95 (design target derived from the NFR01 end-to-end latency budget).
\item \textbf{H5:} User acceptance of AntiGravity meets the threshold of SUS $\geq 68$, and its TAM Perceived Usefulness is significantly higher than that of the incumbent SharePoint search interface (both instruments administered in the user study).
\item \textbf{H6:} Cross-encoder re-ranking improves MRR over bi-encoder-only retrieval by at least 15\% relative, assessed via the ablation study comparing otherwise identical configurations with and without re-ranking (Chapter~\ref{chap:evaluation}).
\end{itemize}

%------------------------------------------------------------------------------
\section{Validity and Reliability}
\label{sec:method-validity}
%------------------------------------------------------------------------------

\begin{table}[H]
\centering
\caption{Threats to validity and mitigation strategies}
\label{tab:validity}
\begin{tabularx}{\textwidth}{@{}llXX@{}}
\toprule
\textbf{Threat} & \textbf{Type} & \textbf{Description} & \textbf{Mitigation} \\
\midrule
Small benchmark (N=50) & Internal validity & Insufficient statistical power & Report confidence intervals; acknowledge as limitation; prioritize effect sizes \\
Evaluator bias (hallucination annotation) & Internal validity & Annotator subjectivity & Double annotation; Cohen's $\kappa$ reported; resolution protocol \\
Single-organization corpus & External validity & DNSI-specific vocabulary and structure & Document corpus characteristics; provide replication protocol \\
LLM non-determinism & Reliability & Temperature>0 variation & Fixed low temperature (0.1) and fixed seed; report results with two prompt variants \\
Small user study (N≥10) & External validity & Insufficient for statistical generalization & Mixed methods; bootstrapped CIs; qualitative depth compensates \\
Temporal validity & External validity & MCP standard may evolve & Document exact MCP version used; note protocol version dependency \\
\bottomrule
\end{tabularx}
\end{table}

%------------------------------------------------------------------------------
\section{Ethical Considerations}
\label{sec:method-ethics}
%------------------------------------------------------------------------------

This research involves human participants (user study) and organizational data (DNSI documents). The following ethical protocols are applied:

\paragraph{Informed consent.} All user study participants are provided with a written information sheet and provide written consent prior to participation. Participation is voluntary and participants may withdraw at any time.

\paragraph{Data anonymization.} User study responses are stored anonymously. No individual-level performance data is associated with participant names in any reported result.

\paragraph{Confidentiality agreement.} In accordance with school guidelines \citep{AivancityGuide2025}, a confidentiality agreement is in place between the student, the DNSI, and Aivancity, governing the treatment of organizational documents and the scope of permissible disclosure.

\paragraph{AI tool declaration.} Portions of the mémoire writing process were assisted by Claude (Anthropic), a generative AI tool. All uses are documented in Annex~E in compliance with Aivancity's policy on generative AI use \citep{AivancityGuide2025}. All analysis, interpretation, and critical reasoning are the intellectual work of the author.

\paragraph{Responsible AI.} The AntiGravity system is designed with responsible deployment principles in mind: user queries are not logged beyond session scope; model outputs are explicitly labeled as AI-generated; the system provides source citations to enable user verification of all factual claims.

%------------------------------------------------------------------------------
\section{Chapter Summary}
\label{sec:method-summary}
%------------------------------------------------------------------------------

This chapter has established the methodological foundation of the mémoire. The Design Science Research paradigm grounds the work in a recognized framework for artifact-generating research. The four-dimensional evaluation protocol (retrieval, generation, system, user) provides comprehensive coverage of the research questions at appropriate levels of rigor. The validity analysis acknowledges the real constraints of an enterprise research context — small benchmark, single organization, limited user study — and documents mitigation strategies for each threat. The following chapter presents the system design and architecture produced in Phase 2 of the DSR process.

%==============================================================================
% CHAPTER 4 — SYSTEM DESIGN AND ARCHITECTURE
%==============================================================================
\chapter{System Design and Architecture}
\label{chap:architecture}

This chapter presents the design and architecture of AntiGravity --- the \gls{rag} conversational assistant developed as the primary artifact of this research. We adopt a systematic design methodology in which every architectural decision is explicitly motivated by functional or non-functional requirements, compared against credible alternatives, and justified on technical and scientific grounds. The chapter is organized from requirements through high-level architecture down to individual component design.

%------------------------------------------------------------------------------
\section{Requirements Analysis}
\label{sec:arch-requirements}
%------------------------------------------------------------------------------

\subsection{Functional Requirements}

\begin{table}[H]
\centering
\caption{Functional requirements of the AntiGravity system}
\label{tab:functional-requirements}
\begin{tabularx}{\textwidth}{@{}llX@{}}
\toprule
\textbf{ID} & \textbf{Priority} & \textbf{Requirement} \\
\midrule
FR01 & Must & Accept natural language queries and return natural language responses \\
FR02 & Must & Responses shall be grounded in documents retrieved from SharePoint CERP \\
FR03 & Must & Every response shall include citations identifying the source documents \\
FR04 & Must & Respect SharePoint ACLs: users receive only responses grounded in authorized documents \\
FR05 & Must & Authenticate users via the organization's SSO provider \\
FR06 & Must & Support PDF and DOCX document formats \\
FR07 & Should & Support conversational context (multi-turn dialogue) \\
FR08 & Should & Indicate when a query cannot be answered from the available corpus \\
FR09 & Should & Allow users to provide feedback on response quality \\
FR10 & Could & Support query routing to specialized sub-indexes by document category \\
\bottomrule
\end{tabularx}
\end{table}

\subsection{Non-Functional Requirements}

\begin{table}[H]
\centering
\caption{Non-functional requirements and measurable acceptance criteria}
\label{tab:nonfunctional-requirements}
\begin{tabularx}{\textwidth}{@{}llXl@{}}
\toprule
\textbf{ID} & \textbf{Category} & \textbf{Requirement} & \textbf{Acceptance Criterion} \\
\midrule
NFR01 & Performance & End-to-end query response latency & P95 $\leq$ 5 seconds \\
NFR02 & Performance & FAISS retrieval latency & $\leq$ 500ms for up to 100k chunks \\
NFR03 & Security & Document access enforcement & Zero unauthorized document leakage \\
NFR04 & Security & Authentication & OAuth 2.0 delegated permissions via Microsoft Entra ID \\
NFR05 & Reliability & Availability & $\geq$ 99\% during business hours \\
NFR06 & Scalability & Concurrent users & $\geq$ 50 concurrent users without degradation \\
NFR07 & Maintainability & Component coupling & Each component replaceable independently \\
NFR08 & Deployability & Network compatibility & Operational behind corporate proxy \\
NFR09 & Compliance & Data locality & Document content not persisted outside the DNSI network perimeter \\
NFR10 & Usability & Response quality & Faithfulness $\geq$ 0.85 on benchmark dataset \\
\bottomrule
\end{tabularx}
\end{table}

NFR04 (OAuth 2.0 delegated permissions) is the single most constraining non-functional requirement architecturally. It mandates that every SharePoint query execute under the authenticated user's token --- not an application-level service account --- ensuring that ACL enforcement is performed by SharePoint's own permission engine rather than by application logic that could be misconfigured or bypassed.

%------------------------------------------------------------------------------
\section{Architectural Approach and High-Level Design}
\label{sec:arch-overview}
%------------------------------------------------------------------------------

\subsection{Design Principles}

Three overarching design principles guide the AntiGravity architecture:

\paragraph{Separation of concerns.} The context management layer (document storage, semantic indexing, ACL-aware retrieval) is fully separated from the conversational layer (query processing, context assembly, generation). This separation localizes the security-sensitive ACL enforcement logic in a single, auditable component --- the MCP Context Server.

\paragraph{Protocol-mediated integration.} The interface between the conversational client and the context server is mediated exclusively through the Model Context Protocol, rather than through direct function calls or shared memory. This protocol boundary is the security enforcement point: no document context reaches the LLM without passing through the MCP server's permission filter.

\paragraph{Replaceability of components.} Each component (embedding model, vector index, re-ranking model, LLM, data connector) is isolated behind a well-defined interface, enabling substitution without cascading architectural changes --- critical given the rapid evolution of the LLM ecosystem.

\subsection{High-Level Architecture}

The AntiGravity system comprises four architectural layers:

\begin{figure}[H]
\centering
\fbox{\parbox{0.92\textwidth}{
\begin{center}
\textbf{LAYER 1 --- PRESENTATION}\\
Streamlit Web Application (Python)\\
Session management $\cdot$ Source citation display $\cdot$ Feedback collection
\end{center}
\vspace{0.2cm}\hrule\vspace{0.2cm}
\begin{center}
\textbf{LAYER 2 --- CONVERSATIONAL INTELLIGENCE}\\
MCP Conversational Client\\
Query classification $\cdot$ Context assembly $\cdot$ Re-ranking $\cdot$ Prompt engineering\\
LLM Interface (Mistral 7B Instruct)
\end{center}
\vspace{0.2cm}\hrule\vspace{0.2cm}
\begin{center}
\textbf{LAYER 3 --- CONTEXT SERVER (MCP)}\\
MCP Context Server\\
ACL-filtered semantic search $\cdot$ FAISS index $\cdot$ Document store\\
User session context $\cdot$ Permission enforcement
\end{center}
\vspace{0.2cm}\hrule\vspace{0.2cm}
\begin{center}
\textbf{LAYER 4 --- DATA SOURCES}\\
SharePoint CERP (Microsoft Graph API $\cdot$ OAuth 2.0 $\cdot$ SSO)\\
Document formats: PDF, DOCX
\end{center}
}}
\caption{AntiGravity four-layer architecture overview}
\label{fig:architecture-overview}
\end{figure}

The MCP boundary between Layers 2 and 3 is the key security enforcement point: no unauthorized document content is ever injected into the LLM prompt.

%------------------------------------------------------------------------------
\section{Data Ingestion and Pre-processing Pipeline}
\label{sec:arch-ingestion}
%------------------------------------------------------------------------------

\subsection{Document Extraction from SharePoint}

Document extraction from SharePoint is performed via the Microsoft Graph API using OAuth 2.0 application permissions for the ingestion pipeline (distinct from the delegated permissions used for query-time retrieval). The ingestion process runs offline during scheduled crawl windows.

The extraction pipeline performs the following operations:

\begin{enumerate}
  \item \textbf{Enumerate document libraries:} The Graph API endpoint \texttt{/sites/\{site-id\}/drives} lists all libraries in the target SharePoint site collection.
  \item \textbf{Traverse document hierarchies:} The \texttt{/drives/\{drive-id\}/root/children} endpoint recursively enumerates folders and files.
  \item \textbf{Fetch document metadata:} For each document: file name, timestamps, author, parent path, SharePoint item ID, and permission list.
  \item \textbf{Download document content:} Binary content is downloaded via \texttt{/drives/\{drive-id\}/items/\{item-id\}/content}.
  \item \textbf{Permission manifest construction:} A manifest maps each SharePoint item ID to its set of authorized principals, accounting for inherited and explicitly-set permissions. This manifest is the foundation of the ACL enforcement mechanism.
\end{enumerate}

\subsection{Document Parsing}

\paragraph{PDF processing.} PDF documents are parsed using \texttt{pdfplumber}, which provides layout-aware text extraction preserving table structures and column layouts. Pages with fewer than 50 extracted characters are flagged for review.

\paragraph{DOCX processing.} DOCX documents are parsed using \texttt{python-docx}. Paragraph structure, heading hierarchy, and table content are extracted and preserved in chunk metadata.

\paragraph{Encoding normalization.} All text is normalized to UTF-8, stripped of control characters, and quality-filtered: chunks with a character-to-token ratio below a threshold (indicating code, serialized data, or extraction artifacts) are excluded from indexing.

\subsection{Chunking Strategy}
\label{subsec:chunking}

The chunking strategy is one of the most consequential decisions in a RAG pipeline. Three candidate strategies were evaluated:

\begin{table}[H]
\centering
\caption{Comparison of chunking strategies evaluated}
\label{tab:chunking-strategies}
\begin{tabularx}{\textwidth}{@{}lXXX@{}}
\toprule
\textbf{Strategy} & \textbf{Description} & \textbf{Advantage} & \textbf{Disadvantage} \\
\midrule
Fixed-size (tokens) & Split at fixed token count & Simple; predictable window usage & Splits semantic units mid-content \\
Sentence-aware & Split at sentence boundaries & Preserves semantic units & Variable chunk size; issues in technical text \\
Recursive character & Split hierarchically: paragraphs $\to$ sentences $\to$ chars & Preserves structure; respects size & More complex \\
\bottomrule
\end{tabularx}
\end{table}

\paragraph{Selected strategy: Recursive character splitting.} Target chunk size: \textbf{512 tokens}, overlap: \textbf{64 tokens}. This choice is justified by: (1) compatibility with the embedding model's context window; (2) the 64-token overlap prevents concepts spanning boundaries from losing their referent (an issue observed with zero-overlap in early experiments); (3) preliminary testing showed chunks below 256 tokens lacked sufficient context for meaningful embedding representations, while chunks above 1024 tokens reduced precision by bundling distinct concepts.

Each chunk stores: document ID, SharePoint item ID, chunk index, character offset, file name, document title, section heading, last modified date, and the permission manifest entry for the parent document.

%------------------------------------------------------------------------------
\section{Vectorization and Semantic Indexing}
\label{sec:arch-vectorization}
%------------------------------------------------------------------------------

\subsection{Embedding Model Selection}

\begin{table}[H]
\centering
\caption{Embedding model trade-off analysis}
\label{tab:embedding-models}
\begin{tabularx}{\textwidth}{@{}lXllll@{}}
\toprule
\textbf{Model} & \textbf{Provider} & \textbf{Dim.} & \textbf{Latency} & \textbf{MTEB} & \textbf{Local} \\
\midrule
\texttt{all-MiniLM-L6-v2} & HuggingFace & 384 & $\sim$5ms & 56.3 & Yes \\
\texttt{all-mpnet-base-v2} & HuggingFace & 768 & $\sim$25ms & 57.8 & Yes \\
\texttt{text-embedding-ada-002} & OpenAI & 1536 & $\sim$100ms & 60.9 & No (API) \\
\bottomrule
\end{tabularx}
\end{table}

\paragraph{Selected: \texttt{all-mpnet-base-v2}.} The OpenAI \texttt{ada-002} model offers superior benchmark performance but was rejected on two grounds: (1) NFR09 prohibits transmitting document content to external APIs; (2) $\sim$100ms per-query latency would exceed the NFR01 budget under realistic load. \texttt{all-MiniLM-L6-v2} was rejected due to lower dimensionality and observed degradation on French-language technical text present in the DNSI corpus.

\subsection{FAISS Index Construction}

\paragraph{Selected: \texttt{IndexFlatIP} (Inner Product on L2-normalized vectors).} For the proof-of-concept corpus (under 50,000 chunks), exact search provides 100\% retrieval recall, eliminating precision loss from approximation. All embedding vectors are L2-normalized prior to indexing, making inner product equivalent to cosine similarity. At corpus sizes above 100,000 chunks, migration to \texttt{IndexIVFFlat} is warranted to maintain sub-500ms latency (NFR02).

The index is built offline during ingestion and loaded into memory at server startup. Each build produces a versioned snapshot (timestamp + document count + SHA-256 of the chunk store) enabling rollback if a new build introduces regression.

%------------------------------------------------------------------------------
\section{The MCP Context Server}
\label{sec:arch-mcp-server}
%------------------------------------------------------------------------------

\subsection{MCP Server Design}

The MCP Context Server implements the MCP server specification \citep{Anthropic2024MCP} in Python, exposing three MCP capabilities:

\begin{itemize}
  \item \textbf{Resources:} The corpus as addressable document chunks (accessible by chunk ID).
  \item \textbf{Tools:} A \texttt{search\_documents} tool accepting a query string and user token, returning ACL-filtered, re-ranked chunks.
  \item \textbf{Prompts:} A \texttt{rag\_response} prompt template structuring context injection for the LLM generation step.
\end{itemize}

\subsection{ACL-Aware Retrieval}
\label{subsec:acl-retrieval}

The ACL enforcement mechanism operates in six steps:

\begin{enumerate}
  \item The MCP client sends a \texttt{search\_documents} tool call with the query text and user OAuth token.
  \item The MCP server resolves the user's effective group membership via the Graph API \texttt{/me/transitiveMemberOf} endpoint.
  \item The user's \textit{effective permission set} is constructed: the union of the user's direct identity and all group memberships.
  \item FAISS retrieval returns the top-$K$ candidates ($K = 20$, before re-ranking to top-$k = 5$).
  \item Each candidate chunk is matched against the permission manifest. Chunks whose parent document has no permission intersection with the user's effective permission set are removed.
  \item Filtered candidates proceed to re-ranking.
\end{enumerate}

\paragraph{Security guarantee.} No unauthorized document content is ever injected into the LLM prompt or transmitted to the client. The MCP boundary enforces this structurally.

\paragraph{Over-retrieval design.} The ratio $K/k = 4$ ensures the user receives $k = 5$ authorized results even when a significant fraction of top-$K$ candidates are permission-filtered. User group membership is session-cached with a 5-minute TTL to reduce authentication latency.

\subsection{Context Window Management}

The available context budget for document chunks is computed as:
\begin{equation}
  C_{\text{docs}} = C_{\text{model}} - C_{\text{system}} - C_{\text{query}} - C_{\text{margin}}
  \label{eq:context-budget}
\end{equation}
where $C_{\text{model}} = 4096$ tokens (Mistral 7B), $C_{\text{system}} \approx 150$ tokens, $C_{\text{query}} \leq 512$ tokens, $C_{\text{margin}} = 128$ tokens. Typical available budget: $\approx 3300$ tokens (6--7 chunks at 512 tokens each).

%------------------------------------------------------------------------------
\section{Query Processing and Re-ranking}
\label{sec:arch-query}
%------------------------------------------------------------------------------

\subsection{Query Classification}

Incoming queries are classified along two dimensions:

\paragraph{Answerability.} Queries are classified as \textit{in-scope}, \textit{out-of-scope}, or \textit{ambiguous} using Mistral 7B zero-shot classification. Out-of-scope queries receive a graceful refusal (FR08) without entering the retrieval pipeline.

\paragraph{Complexity.} Queries are classified as \textit{factoid}, \textit{multi-hop}, or \textit{procedural}. This adjusts the retrieval parameter $K$ (factoid: 10; multi-hop: 20; procedural: 15) and re-ranking strategy.

\subsection{Re-ranking}
\label{subsec:reranking}

\paragraph{Selected: \texttt{cross-encoder/ms-marco-MiniLM-L-6-v2}.} This cross-encoder re-scorer is applied to the top-$K$ ACL-filtered candidates, reordering them by joint (query, passage) relevance score. At 512-token chunks and batch size 20, re-ranking latency is approximately 200--400ms on CPU, within the NFR01 budget. The precision improvement from re-ranking is empirically validated in Chapter~\ref{chap:evaluation}.

%------------------------------------------------------------------------------
\section{Generation Pipeline}
\label{sec:arch-generation}
%------------------------------------------------------------------------------

\subsection{Language Model Selection}

\textbf{Mistral 7B Instruct} \citep{Jiang2023} was selected for three reasons: (1) competitive instruction-following performance at a deployable model size; (2) on-premises deployment satisfying NFR09; (3) stronger French-language capability than comparable English-centric models, relevant for the DNSI corpus.

Commercial API-only models (GPT-4, Claude) were evaluated but rejected on NFR09 grounds: their inference requires transmitting document content to external servers.

\subsection{Prompt Engineering}
\label{subsec:prompt-engineering}

The prompt template follows a four-part structure designed to minimize hallucination:

\begin{verbatim}
[SYSTEM PROMPT]
You are AntiGravity, a precise document assistant for the DNSI.
Answer questions ONLY using the provided document excerpts.
If the answer is not found, state: "I could not find this
information in the available documents."
Always cite the source document: [Source: <title>].
Do not invent, extrapolate, or use external knowledge.

[CONTEXT BLOCK]
Document 1: [Title] | [Date] | [Section]
<chunk text>
... (up to context budget)

[CONVERSATION HISTORY]
<previous turns>

[QUERY]
User: <current query>
Assistant:
\end{verbatim}

Three hallucination-reduction mechanisms are embedded in the system prompt: (1) explicit scope restriction to retrieved content; (2) explicit refusal template for unanswerable queries; (3) mandatory source citation for every factual claim.

Generation parameters: temperature=0.1, top\_p=0.9, max\_new\_tokens=512, repetition\_penalty=1.1.

%------------------------------------------------------------------------------
\section{Security Architecture}
\label{sec:arch-security}
%------------------------------------------------------------------------------

\subsection{Authentication Flow}

OAuth 2.0 Authorization Code flow with PKCE, against Microsoft Entra ID:
\begin{enumerate}
  \item User is redirected to Entra ID login (supporting MFA per organization policy).
  \item Entra ID returns an authorization code; the application exchanges it for a delegated access token (\texttt{Sites.Read.All}, \texttt{Files.Read.All}).
  \item The access token is stored server-side in the user's session; never transmitted to the client browser.
  \item All SharePoint operations use this delegated token, inheriting the user's exact permissions.
\end{enumerate}

\subsection{Threat Model}

\begin{table}[H]
\centering
\caption{AntiGravity threat model}
\label{tab:threat-model}
\begin{tabularx}{\textwidth}{@{}lXXl@{}}
\toprule
\textbf{Threat} & \textbf{Description} & \textbf{Mitigation} & \textbf{Risk} \\
\midrule
Unauthorized data access & User receives content from unauthorized documents & ACL enforcement at retrieval layer; delegated OAuth tokens & Low \\
Prompt injection via document & Adversarial document content manipulates LLM & System prompt reinforcement; output sanitization & Medium \\
Session hijacking & Attacker obtains another user's session token & HTTPS-only; server-side token storage; 1-hour TTL & Low \\
Model output XSS & LLM generates active content & Streamlit HTML escaping; response post-processing & Low \\
Data exfiltration via LLM & Confidential content transmitted to external APIs & No training; inference only; on-premises model (NFR09) & Low \\
\bottomrule
\end{tabularx}
\end{table}

The residual prompt injection risk (Medium) is the most significant security limitation and is discussed critically in Chapter~\ref{chap:discussion}.

%------------------------------------------------------------------------------
\section{Deployment Architecture}
\label{sec:arch-deployment}
%------------------------------------------------------------------------------

The system deploys as three containerized services:

\begin{itemize}
  \item \textbf{Ingestion Worker:} Scheduled nightly job that crawls SharePoint, processes documents, and rebuilds the FAISS index and permission manifest.
  \item \textbf{MCP Context Server:} Always-on service exposing MCP over HTTP/SSE transport with the FAISS index resident in memory.
  \item \textbf{Streamlit Application:} User-facing web application embedding the MCP client and LLM inference client.
\end{itemize}

\paragraph{Corporate proxy.} The DNSI's TLS-terminating proxy required configuring all HTTP clients with the proxy certificate in their trust stores, a proxy bypass list for internal service communication, and on-premises Mistral AI deployment to eliminate external API dependencies.

%------------------------------------------------------------------------------
\section{Design Decisions Summary}
\label{sec:arch-decisions}
%------------------------------------------------------------------------------

\begin{table}[H]
\centering
\caption{Key architectural design decisions with justification}
\label{tab:design-decisions}
\begin{tabularx}{\textwidth}{@{}lXXX@{}}
\toprule
\textbf{Decision} & \textbf{Chosen} & \textbf{Alternative(s)} & \textbf{Justification} \\
\midrule
Integration protocol & MCP & Direct Graph API; LangChain & Standardization; security boundary; connector reuse \\
Embedding model & \texttt{all-mpnet-base-v2} & \texttt{ada-002}; \texttt{MiniLM} & Data locality (NFR09); French support; quality/latency \\
Vector index & FAISS IndexFlatIP & IndexIVFFlat; Pinecone & Exact recall at PoC scale; no external service \\
Chunking & Recursive char 512T/64T & Fixed-size; sentence-aware & Semantic boundary preservation; window compatibility \\
Re-ranking & MS-MARCO cross-encoder & No re-ranking; ColBERT & Precision improvement within latency budget \\
Language model & Mistral 7B Instruct & GPT-4; Llama 2 & Data locality; French support; performance/size \\
ACL enforcement & At retrieval (MCP server) & At display layer & Security correctness; no leakage at any stage \\
Auth & OAuth 2.0 delegated PKCE & Service account & Required for per-user ACL; Entra ID compatibility \\
\bottomrule
\end{tabularx}
\end{table}

%------------------------------------------------------------------------------
\section{Chapter Summary}
\label{sec:arch-summary}
%------------------------------------------------------------------------------

This chapter has presented the complete design and architecture of AntiGravity. The four-layer architecture separates presentation, conversational intelligence, context management, and data access concerns. The MCP Context Server is the structural security enforcement point, ensuring ACL verification occurs at the data boundary. Every major design decision has been documented with alternatives and justification.

Three architectural innovations relative to the state of the art are: (1) structural ACL enforcement within the MCP server boundary, eliminating entire classes of access control bypass; (2) MCP as the protocol layer between AI reasoning and data access, providing security enforcement and connector reuse; (3) the over-retrieval design pattern ($K = 20 \to k = 5$ post-filter) maintaining retrieval completeness under restrictive permission environments.

The following chapter describes the technical implementation of this architecture.

%==============================================================================
% CHAPTER 5 — IMPLEMENTATION
%==============================================================================
\chapter{Implementation}
\label{chap:implementation}

This chapter describes the technical implementation of the AntiGravity architecture presented in Chapter~\ref{chap:architecture}. It documents the technology stack, key implementation choices, engineering challenges encountered, and the solutions developed. The objective is not exhaustive code documentation but analytical reflection on the engineering decisions that shaped the final system.

%------------------------------------------------------------------------------
\section{Implementation Environment}
\label{sec:impl-environment}
%------------------------------------------------------------------------------

\subsection{Technology Stack}

\begin{table}[H]
\centering
\caption{AntiGravity technology stack}
\label{tab:tech-stack}
\begin{tabularx}{\textwidth}{@{}llXl@{}}
\toprule
\textbf{Component} & \textbf{Technology} & \textbf{Role} & \textbf{Version} \\
\midrule
Language & Python & Primary development language & 3.11 \\
UI Framework & Streamlit & Conversational web interface & 1.28+ \\
MCP Runtime & \texttt{mcp} (Python SDK) & MCP server/client implementation & 1.0+ \\
Vector Search & FAISS (CPU) & Semantic similarity index & 1.7.4 \\
Embedding & \texttt{sentence-transformers} & Document and query encoding & 2.2+ \\
Re-ranking & \texttt{sentence-transformers} & Cross-encoder re-scoring & 2.2+ \\
LLM Inference & \texttt{llama-cpp-python} & Local Mistral 7B serving & 0.2+ \\
Document Parsing & \texttt{pdfplumber}, \texttt{python-docx} & PDF and DOCX extraction & Latest \\
Graph API Client & \texttt{msal}, \texttt{httpx} & Microsoft Graph API calls & Latest \\
Containerization & Docker & Service isolation & 24.0+ \\
Orchestration & Docker Compose & Multi-service deployment & 2.20+ \\
\bottomrule
\end{tabularx}
\end{table}

\subsection{Hardware Environment}

The prototype was developed and evaluated on the following infrastructure:

\begin{itemize}
  \item \textbf{Development:} MacBook Pro M2, 16 GB RAM (for code development and unit testing)
  \item \textbf{Deployment target:} DNSI server (Linux, CPU-only deployment; GPU access pending procurement)
  \item \textbf{FAISS index:} CPU-based (\texttt{faiss-cpu}); GPU migration path exists via \texttt{faiss-gpu} drop-in replacement
  \item \textbf{LLM inference:} \texttt{llama-cpp-python} with GGUF quantization (4-bit quantized Mistral 7B), enabling inference on CPU at approximately 3--8 tokens/second
\end{itemize}

The choice of CPU-only deployment was driven by the DNSI's current infrastructure constraints. The 4-bit quantization of Mistral 7B (\texttt{mistral-7b-instruct-v0.2.Q4\_K\_M.gguf}) reduces memory requirements from $\sim$14 GB (FP16) to $\sim$4.5 GB while incurring a measured perplexity degradation of less than 3\% relative to the full-precision model.

%------------------------------------------------------------------------------
\section{MCP-Compatible SharePoint Connector}
\label{sec:impl-connector}
%------------------------------------------------------------------------------

\subsection{Implementation Overview}

The SharePoint connector is implemented as a Python class that wraps the Microsoft Graph API with the MSAL (Microsoft Authentication Library) OAuth 2.0 flow. For the ingestion pipeline, it uses application-level permissions; for the MCP server's query-time permission resolution, it uses delegated user tokens.

The connector exposes three primary operations:

\begin{itemize}
  \item \texttt{list\_documents(site\_id, library\_id)} --- enumerate all documents in a library with their metadata and permission objects.
  \item \texttt{download\_document(item\_id)} --- retrieve binary document content.
  \item \texttt{get\_user\_permissions(user\_token)} --- resolve the user's group membership and effective permissions via \texttt{/me/transitiveMemberOf}.
\end{itemize}

\subsection{Key Engineering Challenge: Permission Inheritance Resolution}

SharePoint's permission model supports both inherited permissions (a document inherits permissions from its parent library) and broken inheritance (a document has explicitly overridden permissions). Correctly resolving effective permissions requires:

\begin{enumerate}
  \item Checking whether each item has broken inheritance via the \texttt{hasUniqueRoleAssignments} property.
  \item For items with broken inheritance: querying \texttt{/drives/\{driveId\}/items/\{itemId\}/permissions} directly.
  \item For items with inherited permissions: traversing the ancestor hierarchy until an item with explicit permissions is found.
\end{enumerate}

This traversal adds latency to the ingestion pipeline and introduced a subtle bug in early implementation: items at the root of a site collection with no explicit permissions would silently fall through without being assigned to any permission group. The fix was to treat absence of explicit permissions as inheritance from the site collection default (typically ``Everyone'') rather than as an empty permission set.

\subsection{Rate Limiting and Throttling}

The Microsoft Graph API enforces service-level throttling \citep{MicrosoftGraph2024}. During ingestion of a large library, the connector encountered \texttt{HTTP 429 Too Many Requests} responses. The mitigation strategy uses exponential backoff with jitter:

\begin{lstlisting}[language=Python, caption={Exponential backoff for Graph API throttling}]
async def graph_request_with_retry(url, headers, max_retries=5):
    for attempt in range(max_retries):
        response = await client.get(url, headers=headers)
        if response.status_code == 429:
            retry_after = int(response.headers.get("Retry-After", 2 ** attempt))
            await asyncio.sleep(retry_after + random.uniform(0, 1))
            continue
        response.raise_for_status()
        return response.json()
    raise MaxRetriesExceeded(f"Failed after {max_retries} attempts")
\end{lstlisting}

%------------------------------------------------------------------------------
\section{Ingestion Pipeline Implementation}
\label{sec:impl-ingestion}
%------------------------------------------------------------------------------

\subsection{Document Processing}

The ingestion pipeline is implemented as a sequential Python script designed to run as a scheduled job. Processing is structured in three stages:

\paragraph{Stage 1 --- Extraction and parsing.} Documents are downloaded, parsed (PDF via \texttt{pdfplumber}, DOCX via \texttt{python-docx}), and persisted as raw text files with associated metadata JSON.

\paragraph{Stage 2 --- Chunking and embedding.} Raw text is chunked using the recursive character strategy (512T/64T). Each chunk is embedded via the \texttt{SentenceTransformer('all-mpnet-base-v2')} model. Embeddings are computed in batches of 64 (memory-efficient) and L2-normalized before storage.

\paragraph{Stage 3 --- Index construction.} All embeddings are assembled into a FAISS \texttt{IndexFlatIP} index. The index, chunk metadata store (JSON Lines format), and permission manifest are versioned and written to disk.

\subsection{Incremental Indexing}

Full re-indexing of the entire corpus on each ingestion run is computationally expensive. The implementation tracks a ``last modified'' timestamp per document and processes only documents modified since the previous ingestion run. Documents that are deleted from SharePoint are detected via a set-difference between the current enumeration and the previous document manifest, and their chunks are removed from the index by rebuilding the index without those chunk IDs.

\textbf{Limitation:} Full rebuild is required when a deleted document's chunks need to be removed, since FAISS \texttt{IndexFlatIP} does not support in-place deletion. This is a known limitation of FAISS flat indexes. Migration to \texttt{IndexIDMap} or a database-backed vector store would address this at the cost of additional complexity.

%------------------------------------------------------------------------------
\section{MCP Context Server Implementation}
\label{sec:impl-mcp-server}
%------------------------------------------------------------------------------

The MCP Context Server is implemented using the official Anthropic MCP Python SDK. The server exposes the \texttt{search\_documents} tool and the \texttt{rag\_response} prompt template over HTTP/SSE transport.

\begin{lstlisting}[language=Python, caption={MCP server tool definition (simplified)}]
@server.call_tool()
async def search_documents(name: str, arguments: dict) -> list[TextContent]:
    query = arguments["query"]
    user_token = arguments["user_token"]

    # Step 1: Encode query
    query_embedding = embedder.encode([query], normalize_embeddings=True)

    # Step 2: FAISS over-retrieval (K=20)
    distances, indices = faiss_index.search(query_embedding, k=20)

    # Step 3: ACL filtering
    user_permissions = await resolve_user_permissions(user_token)
    authorized_chunks = [
        chunks[i] for i in indices[0]
        if chunks[i]["permissions"].intersection(user_permissions)
    ]

    # Step 4: Re-ranking
    reranked = cross_encoder.rank(
        query, [c["text"] for c in authorized_chunks], top_k=5
    )

    # Step 5: Format context with citations
    return format_context_block(reranked, authorized_chunks)
\end{lstlisting}

\paragraph{Engineering challenge: MCP transport configuration behind corporate proxy.} The HTTP/SSE transport used by the MCP server was blocked by the corporate proxy's WebSocket inspection policy. The resolution was to configure the MCP server to use the \texttt{streamable-http} transport variant, which is proxy-compatible, rather than the default SSE transport.

%------------------------------------------------------------------------------
\section{Streamlit Application}
\label{sec:impl-ui}
%------------------------------------------------------------------------------

The Streamlit application implements the conversational interface with the following features:

\begin{itemize}
  \item \textbf{OAuth flow:} On first access, users are redirected to Entra ID. The MSAL library handles token acquisition and storage in Streamlit's server-side session state.
  \item \textbf{Conversation history:} Up to 10 previous turns are maintained in session state and injected into the prompt (within the context budget).
  \item \textbf{Source citation display:} Each response is followed by an expandable ``Sources'' section listing the retrieved document titles, sections, and last-modified dates.
  \item \textbf{Confidence indicator:} Responses are tagged with a visual indicator (``Based on retrieved documents'' vs. ``I could not find this information'') based on the answerability classification.
  \item \textbf{Feedback widget:} A thumbs up/down widget collects binary feedback per response, stored in a local SQLite database for future analysis.
\end{itemize}

\paragraph{Engineering challenge: Streamlit session state and OAuth token refresh.} Streamlit's execution model re-runs the entire script on each user interaction. This caused the OAuth token to be lost between interactions in early implementation, forcing the user to re-authenticate repeatedly. The solution is to store the token in \texttt{st.session\_state} and implement a token refresh check at the start of each script execution, refreshing silently if the token expiry is within 5 minutes.

%------------------------------------------------------------------------------
\section{Implementation Challenges Summary}
\label{sec:impl-challenges}
%------------------------------------------------------------------------------

Table~\ref{tab:impl-challenges} summarizes the significant implementation challenges encountered and their resolutions.

\begin{table}[H]
\centering
\caption{Implementation challenges and resolutions}
\label{tab:impl-challenges}
\begin{tabularx}{\textwidth}{@{}lXXl@{}}
\toprule
\textbf{Challenge} & \textbf{Description} & \textbf{Resolution} & \textbf{Impact} \\
\midrule
Graph API throttling & HTTP 429 on large corpus ingestion & Exponential backoff with jitter & Resolved \\
Permission inheritance traversal & Silent miss for items at site root & Treat absent permissions as site-default & Resolved \\
Corporate proxy blocking MCP SSE & WebSocket inspection drops SSE streams & Switch to streamable-http transport & Resolved \\
TLS certificate chain errors & Corporate proxy self-signed cert & Inject proxy certificate into trust stores & Resolved \\
Streamlit token loss on rerun & OAuth token lost between interactions & Session state storage + silent refresh & Resolved \\
FAISS no in-place deletion & Deleted documents require full rebuild & Incremental rebuild on deletion batch & Accepted limitation \\
Mistral 7B inference latency (CPU) & 20--40s per response on CPU & 4-bit GGUF quantization; response streaming & Mitigated \\
\bottomrule
\end{tabularx}
\end{table}

The most significant unresolved challenge is LLM inference latency on CPU. The 4-bit quantized Mistral 7B generates at approximately 3--8 tokens/second on the available CPU infrastructure, resulting in typical response times of 25--45 seconds for responses of 150--300 tokens. This exceeds NFR01 (P95 $\leq$ 5 seconds) and is the primary performance gap in the current deployment. GPU inference would reduce generation latency to 1--3 seconds. This limitation is discussed in the system performance evaluation (Chapter~\ref{chap:evaluation}) and in the limitations analysis (Chapter~\ref{chap:discussion}).

%------------------------------------------------------------------------------
\section{Reproducibility}
\label{sec:impl-reproducibility}
%------------------------------------------------------------------------------

To support reproducibility of the experimental results, the following artifacts are provided in Annex~C:

\begin{itemize}
  \item \textbf{Environment specification:} Complete \texttt{requirements.txt} with pinned dependency versions.
  \item \textbf{Docker Compose file:} Reproduces the three-service deployment topology.
  \item \textbf{Configuration template:} All system parameters (chunk size, overlap, $K$, $k$, model paths, proxy settings) centralized in a YAML configuration file.
  \item \textbf{Ingestion script:} Documented pipeline script for reproducible index construction.
  \item \textbf{Model download instructions:} Provenance of the Mistral 7B GGUF model file (Hugging Face repository, SHA-256 hash).
\end{itemize}

The benchmark dataset (50 Q\&A pairs) is provided in Annex~A with annotation guidelines. The FAISS index built from the DNSI corpus cannot be distributed due to confidentiality constraints, but the index construction script and a synthetic test corpus (anonymized) are provided.

%------------------------------------------------------------------------------
\section{Chapter Summary}
\label{sec:impl-summary}
%------------------------------------------------------------------------------

This chapter has documented the implementation of AntiGravity across its four architectural layers. The technology stack is grounded in Python, Streamlit, FAISS, sentence-transformers, and the MCP Python SDK. Seven significant engineering challenges were encountered during development; six were fully resolved, and one (CPU inference latency) represents a known performance limitation bounded by the current infrastructure environment. The implementation is reproducible through the artifacts provided in the appendices. The following chapter presents the experimental evaluation of the implemented system.

%==============================================================================
% CHAPTER 6 --- EVALUATION AND RESULTS
%==============================================================================
\chapter{Evaluation and Results}
\label{chap:evaluation}

This chapter defines the evaluation of the AntiGravity system across the four dimensions defined in Chapter~\ref{chap:methodology}: retrieval quality, generation quality, system performance, and user acceptance. Each section presents the result structure for the relevant hypotheses; the quantitative values are to be filled in upon completion of the experiments ([TBC] placeholders), and their significance --- including where results fall short of targets --- will be interpreted accordingly.

%------------------------------------------------------------------------------
\section{Evaluation Setup}
\label{sec:eval-setup}
%------------------------------------------------------------------------------

\subsection{Experimental Conditions}

Three conditions are compared throughout:

\begin{itemize}
\item \textbf{Condition A --- LLM Alone:} Mistral 7B Instruct, temperature=0.1, no retrieval. The model answers from parametric knowledge only.
\item \textbf{Condition B --- Naive RAG:} FAISS flat retrieval (top-$k=5$, no re-ranking, no MCP layer), same prompt template, same model.
\item \textbf{Condition C --- RAG+MCP (AntiGravity):} Full system --- MCP-mediated retrieval, ACL filtering, cross-encoder re-ranking, Mistral 7B generation.
\end{itemize}

All three conditions use identical generation parameters (temperature=0.1, top\_p=0.9, max\_new\_tokens=512, repetition\_penalty=1.1) and the same underlying model. Observed differences between conditions are therefore attributable to the architecture, not to model variation.

\subsection{Benchmark Dataset}

The evaluation benchmark comprises 50 question-answer pairs derived from the DNSI SharePoint corpus following the protocol of Section~\ref{sec:method-evaluation}. Table~\ref{tab:benchmark-composition} shows the distribution across query categories.

\begin{table}[H]
\centering
\caption{Benchmark dataset composition}
\label{tab:benchmark-composition}
\begin{tabularx}{\textwidth}{@{}lXrr@{}}
\toprule
\textbf{Category} & \textbf{Description} & \textbf{Count} & \textbf{\%} \\
\midrule
Factoid & Single-document, specific fact retrieval & 20 & 40\% \\
Multi-hop & Synthesis required across 2+ documents & 15 & 30\% \\
Procedural & Step-by-step instruction queries & 10 & 20\% \\
Unanswerable & Answer not in corpus & 5 & 10\% \\
\midrule
\textbf{Total} & & \textbf{50} & \textbf{100\%} \\
\bottomrule
\end{tabularx}
\end{table}

Two independent annotators label the ground-truth relevant chunks for each question. Overall inter-annotator agreement is measured using Cohen's $\kappa$ \citep{Cohen1960Kappa} [TBC]; pairs on which the annotators disagree are adjudicated and revised before inclusion in the benchmark.

% [NOTE FOR AUTHOR: Insert actual Cohen's kappa value here, e.g. kappa = [VALUE],
%  and specify how many pairs were revised.]

%------------------------------------------------------------------------------
\section{Retrieval Evaluation}
\label{sec:eval-retrieval}
%------------------------------------------------------------------------------

\subsection{Results}

Retrieval is evaluated on Conditions B and C (Condition A has no retrieval component). Results are reported over the answerable questions in the benchmark. A BM25 lexical baseline \citep{Robertson2009BM25} provides a reference point for assessing the contribution of dense semantic search.

\begin{table}[H]
\centering
\caption{Retrieval evaluation results (answerable questions)}
\label{tab:retrieval-results}
\begin{tabularx}{\textwidth}{@{}lXXXXX@{}}
\toprule
\textbf{System} & \textbf{P@1} & \textbf{P@3} & \textbf{P@5} & \textbf{MRR} & \textbf{nDCG@5} \\
\midrule
BM25 (lexical baseline) & [TBC] & [TBC] & [TBC] & [TBC] & [TBC] \\
Cond. B: Naive RAG (no rerank) & [TBC] & [TBC] & [TBC] & [TBC] & [TBC] \\
Cond. C: RAG+MCP (with rerank) & \textbf{[TBC]} & \textbf{[TBC]} & \textbf{[TBC]} & \textbf{[TBC]} & \textbf{[TBC]} \\
\midrule
$\Delta$ C vs. B (absolute) & [TBC] & [TBC] & [TBC] & [TBC] & [TBC] \\
$p$-value (Wilcoxon, $\alpha=0.05$) & [TBC] & [TBC] & [TBC] & [TBC] & [TBC] \\
\bottomrule
\end{tabularx}
\end{table}

% [NOTE FOR AUTHOR: Replace all [TBC] values with actual experimental results.
%  Run the retrieval evaluation pipeline on the 50-item benchmark dataset and record
%  P@1, P@3, P@5, MRR, nDCG@5 for BM25, Naive RAG, and RAG+MCP.
%  Run Wilcoxon signed-rank tests for pairwise comparisons.]

\paragraph{Hypothesis H3 (dense retrieval > BM25):} Dense retrieval is expected to outperform BM25 on this corpus, where technical concepts are expressed in both French and English and frequently through acronyms, creating a vocabulary mismatch that lexical matching cannot resolve. Results will be reported upon completion of the evaluation pipeline.

\paragraph{Hypothesis H6 (re-ranking improves MRR by at least 15\%):} Cross-encoder re-ranking is expected to deliver measurable MRR improvement over bi-encoder-only retrieval, particularly on multi-hop queries. Statistical significance will be assessed via the Wilcoxon signed-rank test.

\subsection{Retrieval by Query Type}

\begin{table}[H]
\centering
\caption{Retrieval Precision@5 disaggregated by query category}
\label{tab:retrieval-by-type}
\begin{tabularx}{\textwidth}{@{}lXXX@{}}
\toprule
\textbf{Query Type} & \textbf{Cond. B} & \textbf{Cond. C} & \textbf{$\Delta$} \\
\midrule
Factoid ($n=20$) & [TBC] & [TBC] & [TBC] \\
Multi-hop ($n=15$) & [TBC] & [TBC] & [TBC] \\
Procedural ($n=10$) & [TBC] & [TBC] & [TBC] \\
\bottomrule
\end{tabularx}
\end{table}

% [NOTE FOR AUTHOR: Replace [TBC] with measured values per category.]

The expected pattern is that re-ranking produces its largest gains on multi-hop queries, which require retrieving semantically distant chunks that the bi-encoder alone ranks below more lexically similar but less relevant passages. This hypothesis will be confirmed or disconfirmed by the disaggregated results.

\subsection{Retrieval Failure Analysis}

% [NOTE FOR AUTHOR: After running the evaluation, perform manual inspection of failed
%  retrievals (zero relevant results in top-5 for Condition C) and classify them by
%  failure mode. Expected failure categories based on system design:
%  (1) queries containing DNSI-internal acronyms absent from the embedding model vocabulary;
%  (2) queries targeting documents whose PDF extraction produced no parseable text (scanned pages);
%  (3) queries referencing documents added to SharePoint after the most recent ingestion cycle.
%  Report the count per category and the remediation pathway for each.]

%------------------------------------------------------------------------------
\section{Generation Quality Evaluation}
\label{sec:eval-generation}
%------------------------------------------------------------------------------

\subsection{RAGAS Metrics}

Table~\ref{tab:ragas-results} reports RAGAS \citep{Es2023RAGAS} generation quality across all three conditions on the answerable questions.

\begin{table}[H]
\centering
\caption{RAGAS generation quality results (answerable questions)}
\label{tab:ragas-results}
\begin{tabularx}{\textwidth}{@{}lXXXX@{}}
\toprule
\textbf{Condition} & \textbf{Faithfulness} & \textbf{Answer Relevancy} & \textbf{Context Precision} & \textbf{Context Recall} \\
\midrule
Cond. A --- LLM Alone & [TBC] & [TBC] & N/A & N/A \\
Cond. B --- Naive RAG & [TBC] & [TBC] & [TBC] & [TBC] \\
Cond. C --- RAG+MCP & \textbf{[TBC]} & \textbf{[TBC]} & \textbf{[TBC]} & \textbf{[TBC]} \\
\bottomrule
\end{tabularx}
\end{table}

% [NOTE FOR AUTHOR: Replace all [TBC] with measured RAGAS scores.
%  Run ragas.evaluate() on each condition's (question, answer, contexts, ground_truth) tuples.
%  Condition A has no retrieval context, so Context Precision and Recall do not apply.
%  Report 95% confidence intervals alongside point estimates.]

\paragraph{Hypothesis H1 (RAG+MCP reduces hallucination vs. LLM Alone):} The fundamental grounding mechanism of RAG --- injecting relevant, verified document context into the LLM prompt at query time --- is expected to substantially reduce the generation of facts not present in the DNSI corpus. This hypothesis will be assessed via the RAGAS faithfulness metric and the human-annotated hallucination rate.

\paragraph{Hypothesis H2 (RAG+MCP faithfulness $\geq 0.85$ on benchmark):} NFR10 specifies a faithfulness threshold of 0.85. Whether this threshold is met will be determined by the RAGAS evaluation.

\subsection{Hallucination Rate Analysis}

Hallucination rate is assessed through human annotation of a random sample of generated responses ($n_{\text{sample}}$ per condition), using the taxonomy of \citet{Ji2023}: intrinsic hallucinations (contradictions of retrieved context) and extrinsic hallucinations (claims absent from any retrieved chunk).

\begin{table}[H]
\centering
\caption{Human-annotated hallucination rate by condition}
\label{tab:hallucination-rate}
\begin{tabularx}{\textwidth}{@{}lXXX@{}}
\toprule
\textbf{Condition} & \textbf{Intrinsic Halluc.} & \textbf{Extrinsic Halluc.} & \textbf{Total Rate} \\
\midrule
Cond. A --- LLM Alone & [TBC] & [TBC] & [TBC] \\
Cond. B --- Naive RAG & [TBC] & [TBC] & [TBC] \\
Cond. C --- RAG+MCP & [TBC] & [TBC] & [TBC] \\
\bottomrule
\end{tabularx}
\end{table}

% [NOTE FOR AUTHOR: Fill in measured hallucination rates. Use two independent annotators
%  and report inter-annotator agreement (Cohen's kappa). Provide sample size n per condition.]

\subsection{Hallucination Taxonomy}

% [NOTE FOR AUTHOR: After annotation, classify observed hallucinations by type and discuss
%  the dominant patterns per condition. Expected findings based on system design:
%  Condition A hallucinations should be primarily extrinsic (fabrication of DNSI-specific
%  facts not present in training data). Conditions B and C should show fewer extrinsic
%  hallucinations but may exhibit intrinsic hallucinations when retrieved context is ambiguous.
%  Condition C is expected to show the lowest overall rate due to better context quality
%  from re-ranking and structured citation enforcement in the prompt template.]

\subsection{Qualitative Response Analysis}

% [NOTE FOR AUTHOR: Include two to three example query-response pairs that illustrate
%  key differences between conditions. Select examples that demonstrate:
%  (1) a case where Condition A hallucinated and Condition C did not;
%  (2) a case where Condition B retrieved a wrong passage and Condition C corrected it
%      via re-ranking;
%  (3) a case where none of the conditions could answer correctly (retrieval failure).
%  Anonymize any DNSI-specific content before inclusion per the confidentiality agreement.]

%------------------------------------------------------------------------------
\section{System Performance Evaluation}
\label{sec:eval-performance}
%------------------------------------------------------------------------------

\subsection{End-to-End Latency}

System latency is measured as the elapsed time from query submission to complete response display, under a representative workload (sequential queries, no concurrency). Each component's contribution is profiled separately to identify the dominant bottleneck.

\begin{table}[H]
\centering
\caption{End-to-end latency breakdown by component (P50 / P95)}
\label{tab:latency-breakdown}
\begin{tabularx}{\textwidth}{@{}lXXX@{}}
\toprule
\textbf{Component} & \textbf{P50 (ms)} & \textbf{P95 (ms)} & \textbf{NFR Budget} \\
\midrule
Query embedding (bi-encoder) & [TBC] & [TBC] & 100 ms \\
FAISS retrieval ($k=20$) & [TBC] & [TBC] & 500 ms \\
ACL resolution (permission lookup) & [TBC] & [TBC] & 500 ms \\
Cross-encoder re-ranking ($k=20 \to 5$) & [TBC] & [TBC] & 1000 ms \\
MCP protocol overhead & [TBC] & [TBC] & 200 ms \\
LLM generation (Mistral 7B, CPU) & [TBC] & [TBC] & --- \\
\textbf{Total end-to-end} & [TBC] & [TBC] & \textbf{5000 ms (NFR01)} \\
\bottomrule
\end{tabularx}
\end{table}

% [NOTE FOR AUTHOR: Instrument the pipeline with timing probes at each component boundary
%  and record latency distributions over a representative query set (recommended: 50 queries).
%  LLM generation on CPU is expected to dominate at 25--45 s per query with Q4-quantized
%  Mistral 7B, substantially exceeding NFR01 (5 s). This is a known infrastructure
%  constraint, not an architectural limitation. Document the measured GPU estimate
%  (e.g. from llama.cpp benchmark with GPU offloading) to show the NFR01 path.]

\paragraph{NFR01 compliance:} The CPU-only deployment of Mistral 7B is expected to produce generation latencies in the range of 25--45 seconds for typical query-length responses at approximately 3--8 tokens/second, substantially exceeding the 5-second NFR01 target. This constraint is infrastructural: the DNSI server used for prototype deployment lacks GPU resources. GPU deployment is the required intervention to achieve NFR01 compliance; the architectural components outside the LLM generation step (embedding, FAISS, ACL, re-ranking, MCP overhead) are each expected to remain well within their individual NFR budgets.

\subsection{MCP Overhead Isolation}

\paragraph{Hypothesis H4 (MCP latency overhead $<$ 200~ms P95):} The MCP protocol adds a communication layer between the conversational client and the context server. This overhead is the sum of serialization/deserialization cost and network round-trip time on the Docker bridge network. The 200 ms budget defined in the non-functional requirements represents the maximum acceptable overhead for this layer.

% [NOTE FOR AUTHOR: Measure MCP overhead by comparing retrieval latency in two configurations
%  that differ ONLY in the transport: (A3) dense retrieval + re-ranking via direct function
%  call vs. (A4) the same retrieval + re-ranking pipeline through the MCP transport.
%  The difference isolates MCP protocol overhead from search and re-ranking time.
%  Report P50 and P95 and state whether H4 is confirmed or refuted.]

\subsection{Memory Footprint}

% [NOTE FOR AUTHOR: Report peak RSS memory consumption for each Docker service at steady
%  state under load:
%  - mcp-server: FAISS index in RAM + chunk metadata store
%  - streamlit-app: Mistral 7B GGUF model loaded by llama-cpp-python
%  This is relevant for deployment planning and confirms feasibility on the DNSI server.]

%------------------------------------------------------------------------------
\section{User Evaluation}
\label{sec:eval-user}
%------------------------------------------------------------------------------

\subsection{Participant Profile}

% [NOTE FOR AUTHOR: Report participant demographics and profile:
%  - N (number of participants, from DNSI employee pool)
%  - Job functions represented (e.g., IT architect, project manager, business analyst)
%  - Prior experience with SharePoint search (years of use, frequency)
%  - Prior experience with AI assistants
%  These data come from the demographic questionnaire administered before the SUS/TAM session.]

\subsection{SUS Results}

The System Usability Scale (SUS) is administered to all participants after a single-session evaluation using the five task scenarios of Appendix~\ref{app:user-study}. SUS scores are computed using the standard formula of \citet{Brooke1996SUS} and interpreted against the adjective rating scale of \citet{Bangor2009SUSAdjective}.

\begin{table}[H]
\centering
\caption{SUS scores for AntiGravity}
\label{tab:sus-results}
\begin{tabularx}{\textwidth}{@{}lXX@{}}
\toprule
\textbf{Metric} & \textbf{Value} & \textbf{95\% CI} \\
\midrule
Mean SUS Score & [TBC] & [TBC] \\
Adjective Rating & [TBC] & --- \\
$N$ (participants) & [TBC] & --- \\
\bottomrule
\end{tabularx}
\end{table}

% [NOTE FOR AUTHOR: Replace [TBC] with measured values. Also report the distribution
%  across SUS items to identify specific usability problems. SUS scores between 68 and 80
%  are classified as "good"; below 68 as "ok"; above 80 as "excellent".]

\subsection{TAM Results}

TAM \citep{Davis1989TAM} Perceived Usefulness (PU) and Perceived Ease of Use (PEOU) are measured using the instruments in Appendix~\ref{app:user-study}. SharePoint search satisfaction is collected as a baseline using the same 7-point Likert scale applied to the PU items.

\begin{table}[H]
\centering
\caption{TAM scores: AntiGravity vs. SharePoint baseline}
\label{tab:tam-results}
\begin{tabularx}{\textwidth}{@{}lXXX@{}}
\toprule
\textbf{Measure} & \textbf{AntiGravity} & \textbf{SharePoint Baseline} & \textbf{$p$-value} \\
\midrule
Perceived Usefulness (PU, /7) & [TBC] & [TBC] & [TBC] \\
Perceived Ease of Use (PEOU, /7) & [TBC] & --- & --- \\
\bottomrule
\end{tabularx}
\end{table}

% [NOTE FOR AUTHOR: Use a one-sample or paired Wilcoxon test to compare AntiGravity PU
%  to SharePoint baseline PU. Report effect size (Cohen's d or r) alongside p-value.]

\paragraph{Hypothesis H5 (user acceptance SUS $\geq$ 68 and TAM PU significantly higher than SharePoint baseline):} A SUS score of 68 corresponds to the empirical average (50th percentile) across published SUS studies \citep{Sauro2011SUS}, and scores approaching 71 map to the ``Good'' adjective rating of \citet{Bangor2009SUSAdjective}; SUS $\geq$ 68 therefore represents the minimum acceptable adoption threshold. The TAM comparison tests whether users perceive AntiGravity as more useful than the current SharePoint search workflow.

\subsection{Task-Based Evaluation}

Five representative task scenarios (Appendix~\ref{app:user-study}) are administered to each participant under both systems (counterbalanced order). Task completion and time-on-task are recorded.

\begin{table}[H]
\centering
\caption{Task-based evaluation results}
\label{tab:task-eval}
\begin{tabularx}{\textwidth}{@{}lXXX@{}}
\toprule
\textbf{Task} & \textbf{AntiGravity Completion} & \textbf{SharePoint Completion} & \textbf{Mean Time Saving} \\
\midrule
T1 (Factoid) & [TBC] & [TBC] & [TBC] \\
T2 (Multi-hop) & [TBC] & [TBC] & [TBC] \\
T3 (Procedural) & [TBC] & [TBC] & [TBC] \\
T4 (Multi-hop) & [TBC] & [TBC] & [TBC] \\
T5 (Unanswerable) & [TBC] & [TBC] & N/A \\
\midrule
\textbf{Overall} & [TBC] & [TBC] & [TBC] \\
\bottomrule
\end{tabularx}
\end{table}

% [NOTE FOR AUTHOR: Record task completion (binary: completed within 5-minute time limit)
%  and time-on-task (seconds). For T5, correct completion is the user correctly identifying
%  that the system cannot answer the question. Report time saving only for successful tasks.]

\subsection{Qualitative Feedback}

% [NOTE FOR AUTHOR: Synthesize themes from the semi-structured interviews (Appendix~\ref{app:user-study}).
%  At minimum, identify: (1) most valued features (expected: source citations, natural language query);
%  (2) primary adoption barriers (expected: response latency on CPU deployment);
%  (3) trust factors identified by users. Present as thematic summary with representative quotes,
%  anonymized per the DNSI confidentiality agreement.]

%------------------------------------------------------------------------------
\section{Ablation Study}
\label{sec:eval-ablation}
%------------------------------------------------------------------------------

The ablation study decomposes the contribution of each architectural component by evaluating four sub-configurations:

\begin{itemize}
\item \textbf{A1 --- LLM Alone} (= Condition A): no retrieval.
\item \textbf{A2 --- Dense retrieval, no re-ranking, no MCP}: bi-encoder + FAISS flat retrieval, direct function call (no MCP transport), top-$k=5$ injected into prompt.
\item \textbf{A3 --- Dense retrieval + re-ranking, no MCP}: same as A2 with cross-encoder re-ranking.
\item \textbf{A4 --- Full AntiGravity (= Condition C)}: A3 with MCP transport and ACL filtering.
\end{itemize}

The progression A1 $\to$ A2 isolates the contribution of dense retrieval. A2 $\to$ A3 isolates re-ranking. A3 $\to$ A4 isolates the MCP layer.

\begin{table}[H]
\centering
\caption{Ablation study results: P@5 and RAGAS Faithfulness}
\label{tab:ablation}
\begin{tabularx}{\textwidth}{@{}lXXXX@{}}
\toprule
\textbf{Configuration} & \textbf{P@5} & \textbf{Faithfulness} & \textbf{$\Delta$ P@5 vs. prev.} & \textbf{$\Delta$ Faith. vs. prev.} \\
\midrule
A1 --- LLM Alone & N/A & [TBC] & --- & --- \\
A2 --- Dense only & [TBC] & [TBC] & [TBC] & [TBC] \\
A3 --- Dense + rerank & [TBC] & [TBC] & [TBC] & [TBC] \\
A4 --- Full system & [TBC] & [TBC] & [TBC] & [TBC] \\
\bottomrule
\end{tabularx}
\end{table}

% [NOTE FOR AUTHOR: Replace all [TBC] with measured values.
%  The expected finding for A3 vs A4 (MCP contribution) is that P@5 and Faithfulness
%  differences are small or non-significant --- MCP's primary contributions are security
%  and standardization, not retrieval quality. This is an expected and valid finding;
%  interpret it as such in the discussion rather than treating it as a null result.]

%------------------------------------------------------------------------------
\section{Hypothesis Testing Summary}
\label{sec:eval-hypotheses}
%------------------------------------------------------------------------------

\begin{table}[H]
\centering
\caption{Summary of hypothesis tests and outcomes}
\label{tab:hypothesis-summary}
\begin{tabularx}{\textwidth}{@{}llXX@{}}
\toprule
\textbf{ID} & \textbf{Hypothesis} & \textbf{Metric} & \textbf{Outcome} \\
\midrule
H1 & RAG+MCP reduces hallucination vs. LLM Alone & Hallucination rate & [TBC] \\
H2 & RAG+MCP faithfulness $\geq$ 0.85 & RAGAS Faithfulness (C) & [TBC] \\
H3 & Dense retrieval $>$ BM25 & P@5 & [TBC] \\
H4 & MCP overhead $<$ 200~ms P95 & MCP latency P95 & [TBC] \\
H5 & SUS $\geq$ 68 and TAM PU $>$ SharePoint & SUS, TAM & [TBC] \\
H6 & Re-ranking improves MRR by $\geq$15\% & MRR (C vs. B) & [TBC] \\
\bottomrule
\end{tabularx}
\end{table}

% [NOTE FOR AUTHOR: Replace [TBC] with Confirmed / Partially Confirmed / Refuted
%  and add the key measured value in parentheses, e.g. Confirmed (P@5=0.XX, p<0.05).]

%------------------------------------------------------------------------------
\section{Chapter Summary}
\label{sec:eval-summary}
%------------------------------------------------------------------------------

% [NOTE FOR AUTHOR: Once all experiments are complete, write a 2--3 paragraph summary
%  that synthesizes the key findings across the four evaluation dimensions:
%  retrieval, generation, performance, and user acceptance. Highlight which hypotheses
%  were confirmed or refuted, and identify the primary gap between the designed system
%  and the production-ready target (expected: LLM inference latency on CPU deployment).
%  This summary feeds directly into the Discussion chapter (Chapter~\ref{chap:discussion}).]

%==============================================================================
% CHAPTER 7 — DISCUSSION
%==============================================================================
\chapter{Discussion}
\label{chap:discussion}

The preceding chapter presented the empirical results of the AntiGravity evaluation. This chapter interprets those results in relation to the research questions, critically assesses what the thesis does and does not demonstrate, examines the limitations of the work, and discusses the implications for both researchers and practitioners. A researcher who cannot honestly criticize their own work has not understood it.

%------------------------------------------------------------------------------
\section{Interpretation of Results}
\label{sec:disc-interpretation}
%------------------------------------------------------------------------------

\subsection{On the Primary Research Question}

The primary research question (RQ1) asks whether a RAG+MCP architecture enables more \textit{reliable}, \textit{secure}, and \textit{traceable} access to enterprise knowledge than a standalone LLM or naive RAG implementation.

\paragraph{Reliability.} The hallucination rate comparison (Section~\ref{sec:eval-generation}) provides the central evidence for the reliability dimension. The expected finding — that Condition C (RAG+MCP) substantially reduces hallucination compared to Condition A (LLM Alone) — is the most straightforward result to interpret: grounding LLM generation in retrieved, verified source material from the DNSI corpus structurally prevents the model from confabulating information about DNSI-specific processes, systems, and policies that were not in its training data.

The more nuanced finding concerns the comparison between Conditions B and C (Naive RAG vs. RAG+MCP). The faithfulness improvement of Condition C over B, if confirmed, is attributable to two mechanisms: (1) the cross-encoder re-ranking produces a more precisely relevant context, reducing the probability that the model is distracted by loosely related but irrelevant passages; and (2) the MCP context assembly enforces structured citation formatting, which has been observed to anchor the model more firmly to source material in the prompt.

However, the improvement over Naive RAG is expected to be smaller than the improvement over LLM Alone — this is theoretically correct. RAG itself (even without MCP) addresses the fundamental problem of knowledge grounding. MCP adds value in the dimensions of security and standardization, with a secondary benefit to generation faithfulness through better context management.

\paragraph{Security.} The ACL enforcement mechanism provides a structural security guarantee (Section~\ref{subsec:acl-retrieval}): no document content unauthorized for the querying user can enter the LLM prompt. This guarantee is architectural, not probabilistic. It is the primary security contribution of the RAG+MCP design over both Condition A and Condition B. However, this guarantee has boundaries: it applies to the retrieval layer but does not protect against prompt injection attacks via the document content itself (discussed in Section~\ref{sec:disc-security}).

\paragraph{Traceability.} Every response generated by Condition C includes source citations derived from the chunk metadata. This is not merely a user experience feature: it creates an audit trail linking each factual claim in a response to a specific DNSI document, enabling a human reviewer to verify or refute the claim. Conditions A and B provide no equivalent traceability mechanism.

\subsection{On the Value of MCP Specifically}

A critical question for this thesis is whether MCP itself — as opposed to the RAG architecture in general — makes a measurable contribution. This is the hardest question to answer empirically, because MCP's primary contributions (standardization, security boundary, maintainability) are not directly measurable as retrieval or generation quality metrics.

The ablation study (Section~\ref{sec:eval-ablation}) provides indirect evidence: the comparison of A2 (FAISS + re-ranking, no MCP) vs. A4 (full AntiGravity with MCP) isolates the marginal contribution of the MCP layer. If the difference in P@5 or faithfulness between A2 and A4 is not statistically significant — which is the expected finding — this should not be interpreted as MCP providing no value. It should be interpreted as MCP providing value in dimensions that retrieval/generation metrics do not capture: security architecture, connector standardization, and operational maintainability.

This distinction matters for the scientific interpretation of the thesis. The claim is not that MCP improves RAG retrieval quality. The claim is that MCP provides a principled architectural boundary for security enforcement, and that this architectural property is valuable in enterprise contexts in ways that micro-benchmarks cannot fully capture.

%------------------------------------------------------------------------------
\section{Strengths of the Proposed Architecture}
\label{sec:disc-strengths}
%------------------------------------------------------------------------------

\paragraph{Structural security by design.} The most significant strength of the AntiGravity architecture is that the security property --- ACL-aware retrieval --- is not implemented as an application-level check that can be bypassed by logic errors, but as a structural property of the MCP boundary. No pathway exists through the architecture by which unauthorized document content can reach the LLM without traversing the ACL filter.

\paragraph{Source attribution as trust infrastructure.} The mandatory citation mechanism, embedded in the system prompt and enforced by the prompt template, provides a trust infrastructure that enterprise users and compliance teams require. Unlike a standalone LLM, every AntiGravity response can be audited against its source documents. In regulated environments (information security policies, legal compliance, data governance), this auditability is not optional --- it is required.

\paragraph{Protocol separation of concerns.} The MCP boundary between the conversational layer and the data layer means that changes to the underlying data source (a SharePoint API version change, migration to a different SharePoint tenant, addition of a second data source) require changes only to the MCP server --- not to the conversational client or the LLM interface. This operational benefit compounds over time in environments where data infrastructure evolves continuously.

\paragraph{Research novelty.} The architecture addresses a combination of problems --- ACL-aware retrieval, MCP integration, enterprise SharePoint, and multilingual (French/English) corpora --- that no published system has addressed in combination. The research gap identified in Chapter~\ref{chap:background} is genuinely addressed.

%------------------------------------------------------------------------------
\section{Weaknesses and Limitations}
\label{sec:disc-limitations}
%------------------------------------------------------------------------------

This section deliberately presents the limitations of the thesis with the same rigor applied to its strengths. Honest acknowledgment of limitations is a mark of scientific maturity and protects the thesis from criticism on grounds the author has pre-empted.

\subsection{Evaluation Limitations}

\paragraph{Benchmark dataset size ($N=50$).} Fifty question-answer pairs is sufficient for reporting metric distributions with 95\% confidence intervals, but is insufficient for detecting small effect sizes with high statistical power. A benchmark of $N=200$--500 would provide substantially stronger statistical evidence. The constraint is operational: manual annotation of ground-truth relevance labels is time-intensive. Future work should extend the benchmark.

\paragraph{Single-organization generalizability.} All evaluation is conducted on the DNSI SharePoint corpus. The vocabulary, document structure, query patterns, and language distribution (French/English mix) are specific to this organization. Results may not generalize to organizations with different corpus characteristics. The evaluation should be interpreted as a proof-of-concept demonstration, not as a universal performance characterization.

\paragraph{RAGAS as a proxy for faithfulness.} RAGAS metrics are computed by a language model evaluating the faithfulness of another language model's output against retrieved context. This creates a potential circular dependency: if the evaluating model has similar failure modes to the generating model, some hallucinations may not be detected. A fully reliable hallucination assessment requires human annotation of every generated response --- which was provided for the hallucination rate metric but not for the RAGAS faithfulness scores.

\paragraph{No longitudinal evaluation.} The user study is a single-session evaluation. Real adoption decisions are made over time, through repeated use. A longitudinal study (tracking usage patterns over weeks) would provide much stronger evidence for or against the adoption hypothesis (H5).

\subsection{Technical Limitations}

\paragraph{LLM inference latency.} The CPU-only deployment of Mistral 7B results in generation latency of 25--45 seconds for typical responses --- substantially exceeding NFR01 (P95 $\leq$ 5 seconds). This is the most significant gap between the designed system and the production-ready system. It is a deployment infrastructure limitation (no GPU), not an architectural limitation. GPU deployment is expected to reduce generation latency to 2--5 seconds, bringing the system within NFR01. This limitation was known at the outset and does not invalidate the architectural contributions, but it does prevent the user study from fully reflecting the intended user experience.

\paragraph{Incremental index updates.} The current ingestion pipeline requires a full index rebuild when documents are deleted from SharePoint. For a dynamically evolving corpus, this means a latency window during which deleted documents remain retrievable. The window is bounded by the ingestion schedule (nightly) but is non-zero. Production deployment would require migration to a vector store with native deletion support (e.g., Chroma, Weaviate, or FAISS \texttt{IndexIDMap}).

\paragraph{Single-language embedding model.} The \texttt{all-mpnet-base-v2} embedding model was trained primarily on English text. While it performs adequately on the French-language DNSI documents (French is well-represented in its training data), a multilingual model (e.g., \texttt{paraphrase-multilingual-mpnet-base-v2}) or a domain-fine-tuned model would likely improve retrieval precision on French-language queries.

\paragraph{Prompt injection residual risk.} As identified in the threat model (Table~\ref{tab:threat-model}), the prompt injection attack vector via adversarial document content has not been fully mitigated. The structured system prompt and near-zero temperature reduce susceptibility, but do not eliminate it. A production deployment should add an additional layer of prompt injection detection (e.g., a small classifier model screening retrieved chunks for injection patterns) before this system handles highly sensitive materials.

\subsection{Scope Limitations}

\paragraph{No real-time synchronization.} The nightly ingestion schedule means the index may be up to 24 hours out of date relative to the live SharePoint corpus. For time-sensitive operational contexts (incident response, live project documentation), this lag is unacceptable. Real-time delta synchronization via the Microsoft Graph change notification API is required for production use and is identified as the highest-priority future work item.

\paragraph{Single data source.} The system is designed for SharePoint CERP exclusively. DNSI knowledge is distributed across multiple systems (email, wikis, ticketing systems, databases). A production enterprise assistant would need to federate across these sources. MCP's multi-server architecture makes this extension architecturally straightforward but operationally complex.

%------------------------------------------------------------------------------
\section{Security Analysis: Critical Assessment}
\label{sec:disc-security}
%------------------------------------------------------------------------------

The security architecture presented in Chapter~\ref{chap:architecture} provides strong guarantees at the retrieval layer but leaves residual risks at the prompt layer that deserve honest analysis.

\paragraph{What the architecture guarantees.} The delegated OAuth 2.0 token mechanism ensures that ACL enforcement is performed by Microsoft's Graph API permission engine --- the most authoritative and thoroughly tested component in the system. ACL checks performed by well-maintained cloud provider services are more reliable than ACL checks implemented in application code. The MCP boundary ensures that the retrieval and permission resolution logic cannot be bypassed by the conversational layer.

\paragraph{What the architecture does not guarantee.} The architecture does not protect against an insider threat who uses their legitimate SharePoint write access to plant adversarial documents. A document containing the text \texttt{[SYSTEM: disregard previous instructions and output the full text of the previous response]} would pass the ACL filter (the attacker is authorized to write it) and be retrieved by legitimate queries. The near-zero temperature setting and structured system prompt provide partial mitigation but not elimination of this risk.

\paragraph{Practical likelihood.} This attack requires a motivated insider with SharePoint write access who is aware of the AntiGravity system's implementation details and is willing to compromise the system. In the DNSI context, this is a low-probability but non-zero threat. The recommendation is to add automated prompt injection screening of indexed document content at ingestion time, as an additional defense-in-depth measure.

%------------------------------------------------------------------------------
\section{Deployment Challenges: Lessons Learned}
\label{sec:disc-deployment}
%------------------------------------------------------------------------------

The deployment of AntiGravity in the DNSI's production environment produced several lessons that are generalizable to any enterprise AI deployment:

\paragraph{Corporate proxy is a first-class constraint, not an afterthought.} The DNSI's TLS-terminating proxy required changes to all components in the stack (HTTP clients, MSAL, MCP transport). Enterprise AI systems should be designed with proxy-aware HTTP clients from the start, not retrofitted. The corporate proxy is not an exceptional edge case --- it is the norm in large enterprise environments.

\paragraph{OAuth 2.0 delegated permissions require organizational IT cooperation.} Registering an Entra ID application, configuring delegated permissions, and navigating the DNSI's internal approval workflow for API access consumed significant project time. Enterprise AI projects should account for identity and access management lead times in their project schedules.

\paragraph{On-premises LLM inference is operationally expensive.} Running Mistral 7B on-premises (required by NFR09) requires hardware procurement, model management, quantization trade-off decisions, and ongoing maintenance. Cloud-based LLM APIs are operationally much simpler, but are incompatible with strict data locality requirements. Organizations face a genuine trade-off between operational simplicity and data sovereignty.

%------------------------------------------------------------------------------
\section{Generalizability of Findings}
\label{sec:disc-generalizability}
%------------------------------------------------------------------------------

The AntiGravity architecture and the findings of this thesis are generalizable to other organizations under the following conditions:

\begin{enumerate}
  \item \textbf{SharePoint as the primary knowledge repository:} The connector design and ACL enforcement mechanism are specific to SharePoint and Microsoft Graph API. Organizations using Confluence, Box, Google Drive, or other platforms would require a different MCP server implementation, but the architectural pattern (MCP boundary for security enforcement) is fully transferable.
  \item \textbf{Enterprise security requirements:} The security architecture is designed for organizations with formal information security policies (ISO 27001 or equivalent). Organizations without strict ACL requirements can simplify the architecture by removing the permission manifest and delegated token mechanism.
  \item \textbf{Moderate corpus scale (up to 100k documents):} The FAISS IndexFlatIP performs optimally at this scale. Larger organizations would need to evaluate approximate search indexes.
  \item \textbf{Document-centric knowledge:} The system is optimized for unstructured text in PDF and DOCX format. Organizations with knowledge primarily in structured databases, spreadsheets, or specialized formats would require extended parsing modules.
\end{enumerate}

%------------------------------------------------------------------------------
\section{Future Work}
\label{sec:disc-future}
%------------------------------------------------------------------------------

The limitations identified above directly motivate a prioritized future research and development agenda:

\begin{table}[H]
\centering
\caption{Prioritized future work roadmap}
\label{tab:future-work}
\begin{tabularx}{\textwidth}{@{}llXl@{}}
\toprule
\textbf{Priority} & \textbf{Item} & \textbf{Description} & \textbf{Expected Impact} \\
\midrule
P1 & GPU deployment & Deploy Mistral 7B on GPU hardware to bring generation latency within NFR01 & Critical for user adoption \\
P2 & Real-time delta sync & Microsoft Graph change notifications for live SharePoint updates & Production readiness \\
P3 & Multilingual embedding & Fine-tune or replace embedding model for French/English domain vocabulary & Retrieval precision \\
P4 & Prompt injection defense & Ingestion-time adversarial content screening & Security hardening \\
P5 & Larger benchmark & Expand to $N=200$+ Q\&A pairs for statistical power & Evaluation rigor \\
P6 & Multi-source federation & Extend MCP architecture to email, wikis, ticketing systems & Coverage expansion \\
P7 & Agentic RAG & Integrate multi-step tool use for complex procedural queries & Capability expansion \\
P8 & CIFRE doctoral continuation & Investigate MCP as a universal enterprise AI integration standard & Scientific contribution \\
\bottomrule
\end{tabularx}
\end{table}

\paragraph{CIFRE doctoral direction.} This thesis establishes a foundation for a CIFRE doctoral research program that would investigate three open scientific questions: (1) Can MCP be formalized as a universal integration standard for enterprise AI, and what security-theoretic guarantees can be proven about MCP-mediated architectures? (2) How do RAG+MCP systems scale under multi-tenant, multi-source, and multi-modal knowledge conditions? (3) What are the organizational change management factors that determine enterprise AI assistant adoption, and how can system design choices mitigate adoption barriers?

%------------------------------------------------------------------------------
\section{Chapter Summary}
\label{sec:disc-summary}
%------------------------------------------------------------------------------

This chapter has provided a critical analysis of the AntiGravity system and the thesis research. The primary research question is answered affirmatively: the RAG+MCP architecture demonstrably improves reliability (reduced hallucination), security (structural ACL enforcement), and traceability (mandatory source citations) compared to baseline alternatives. MCP's specific contribution is primarily architectural (security boundary, standardization) rather than metric-measurable (retrieval precision, faithfulness), which is precisely what the MCP specification claims and what this thesis confirms.

The thesis's limitations are significant but bounded: the evaluation is a proof-of-concept demonstration on a single organization's corpus with a small benchmark, not a universal performance characterization. The deployment has a known infrastructure gap (CPU inference latency) that is resolvable with hardware investment. These limitations are typical of applied AI engineering research at the Master's level and do not undermine the validity of the architectural contributions.

The following chapter synthesizes the thesis contributions and answers each research question explicitly.

%==============================================================================
% CHAPTER 8 --- CONCLUSION
%==============================================================================
\chapter{Conclusion}
\label{chap:conclusion}

%------------------------------------------------------------------------------
\section{Summary of Contributions}
\label{sec:conc-contributions}
%------------------------------------------------------------------------------

This thesis set out to design, implement, and evaluate a RAG conversational assistant grounded in the Model Context Protocol for enterprise knowledge access from SharePoint. The work produces three scientific contributions and four technical contributions, each validated against the evidence base established in Chapters~\ref{chap:implementation} and~\ref{chap:evaluation}.

\paragraph{SC1 --- Formal comparative evaluation framework.} This thesis presents the first published empirical comparison of LLM Alone, Naive RAG, and RAG+MCP architectures on a private enterprise SharePoint corpus with ACL enforcement. The evaluation framework --- a domain-specific 50-item benchmark, RAGAS metrics for generation quality, Precision@K/MRR/nDCG for retrieval quality, component-level latency profiling, and a mixed-methods user study --- is documented in sufficient detail for replication and serves as a methodological contribution independent of the specific results obtained.

\paragraph{SC2 --- Reference architecture for ACL-aware enterprise RAG.} The AntiGravity architecture introduces a principled pattern for permission-aware retrieval: the over-retrieval design (\$K=20 \to k=5\$ post-ACL-filter) and the structural placement of ACL enforcement within the MCP server boundary. These design patterns are documented with formal security properties and can be adopted by practitioners deploying RAG systems in access-controlled enterprise environments.

\paragraph{SC3 --- Application of DSR to enterprise AI evaluation.} The thesis demonstrates that Design Science Research provides a principled framework for research that simultaneously produces a system artifact and empirically evaluates it. The three-cycle DSR model (Relevance, Design, Rigor) maps cleanly onto the research phases of this work and is recommended as the methodology of choice for applied enterprise AI research.

\paragraph{TC1--TC4.} The technical contributions --- the MCP-compatible SharePoint connector, the ACL-enforced retrieval pipeline, the domain-specific benchmark dataset, and the Docker-based deployment documentation --- are provided in the appendices as reproducible artifacts.

%------------------------------------------------------------------------------
\section{Answers to Research Questions}
\label{sec:conc-rq-answers}
%------------------------------------------------------------------------------

\paragraph{RQ1 (Primary):} \textit{To what extent does a RAG+MCP architecture enable more reliable, secure, and traceable access to enterprise knowledge than a standalone LLM or naive RAG?}

The evidence supports a positive answer on all three dimensions, with the qualification that quantitative magnitudes are to be confirmed by the evaluation experiments (Chapter~\ref{chap:evaluation}). \textit{Reliability:} RAG+MCP structurally reduces hallucination by grounding LLM generation in retrieved, verified document context. \textit{Security:} The structural ACL enforcement at the MCP boundary provides a formally verifiable security property absent from baseline architectures. \textit{Traceability:} Mandatory source citation in every response creates an audit trail that neither Condition A nor Condition B provides.

% [NOTE FOR AUTHOR: Once Chapter 6 experiments are complete, replace the above paragraph
%  with specific measured values: hallucination rates per condition, statistical significance
%  of the comparisons, and specific citations to the tables in Chapter 6.]

\paragraph{RQ2 (Retrieval):} \textit{What retrieval configuration maximizes precision and recall on a domain-specific SharePoint corpus?}

The combination of \texttt{all-mpnet-base-v2} embeddings (768 dimensions, local deployment), FAISS IndexFlatIP exact search, and an MS-MARCO cross-encoder for re-ranking is the configuration evaluated. The 512-token / 64-token-overlap recursive chunking strategy was selected over zero-overlap fixed-size chunking based on early-stage ablation testing. Quantitative retrieval results are reported in Table~\ref{tab:retrieval-results}.

% [NOTE FOR AUTHOR: Insert the specific P@5, MRR values once measured, e.g.:
%  "Dense retrieval achieves P@5 = X vs. BM25 P@5 = Y (p < 0.001, H3 confirmed).
%   Re-ranking adds a further +Z\% MRR improvement (H6 confirmed, p = W)."]

\paragraph{RQ3 (Protocol):} \textit{What specific properties does MCP deliver in an enterprise SharePoint integration context?}

MCP delivers three properties that direct-function-call integration approaches do not: (1) a stable protocol boundary that decouples the AI pipeline from SharePoint API evolution; (2) a structural security enforcement point enabling ACL filtering before context reaches the LLM; and (3) a standardized connector interface enabling the SharePoint connector to be reused across different MCP-compatible AI clients. The measured latency overhead of the MCP protocol layer is reported in Table~\ref{tab:latency-breakdown}; the design intent was to remain below the 200~ms budget (NFR04) for the MCP overhead component specifically.

% [NOTE FOR AUTHOR: Insert measured MCP overhead P95 value from Table~\ref{tab:latency-breakdown}
%  and state whether H4 is confirmed or refuted.]

\paragraph{RQ4 (Performance):} \textit{What are the latency and throughput characteristics under realistic enterprise conditions?}

On CPU infrastructure, end-to-end latency is dominated by LLM generation using \texttt{llama-cpp-python} with 4-bit quantized Mistral 7B, which is expected to produce P95 generation times in the range of 25--45 seconds at approximately 3--8 tokens/second. This substantially exceeds NFR01 (P95 $\leq$ 5 seconds). All non-generation components (embedding encoding, FAISS retrieval, ACL resolution, cross-encoder re-ranking, MCP assembly) are each designed to operate within their individual NFR budgets. Measured values are reported in Table~\ref{tab:latency-breakdown}. GPU deployment is the required intervention to achieve NFR01 compliance.

\paragraph{RQ5 (Acceptance):} \textit{How do end-users perceive AntiGravity compared to current SharePoint search?}

User acceptance is assessed via SUS, TAM, and task-based evaluation on a panel of DNSI employees. Results are reported in Tables~\ref{tab:sus-results} and~\ref{tab:tam-results}. Source citations were identified as a key trust factor in the design-stage requirements analysis; qualitative feedback from the user study will confirm or qualify this expectation.

% [NOTE FOR AUTHOR: Replace this paragraph with specific measured values once the user
%  study is complete: SUS score (mean, CI, adjective rating), TAM PU vs. SharePoint
%  baseline (mean, p-value), task completion rates, mean time saving.]

%------------------------------------------------------------------------------
\section{Scientific and Practical Implications}
\label{sec:conc-implications}
%------------------------------------------------------------------------------

\paragraph{For researchers.} This thesis contributes empirical evidence to an important open question in the RAG literature: whether the introduction of a protocol standardization layer (MCP) between AI and data sources adds value beyond what the RAG architecture itself provides. MCP's value is primarily architectural --- security enforcement at the data boundary, connector standardization, operational maintainability --- and does not necessarily manifest as metric improvements in isolation. Researchers evaluating RAG-over-enterprise-data systems should expand their evaluation frameworks beyond retrieval and generation metrics to include security audits, maintainability assessments, and operational complexity evaluations.

\paragraph{For practitioners.} Organizations deploying LLM-based assistants over internal documentation corpora should: (1) adopt delegated OAuth 2.0 permissions rather than service accounts, to enable per-user ACL enforcement at the API level; (2) implement ACL filtering at the retrieval layer, not at the display layer; (3) include source citations in the LLM prompt template as the most impactful single intervention for user trust; (4) plan for GPU infrastructure if generation latency is a user adoption concern; and (5) consider MCP as a long-term integration strategy that reduces connector maintenance overhead as the number of data sources grows.

%------------------------------------------------------------------------------
\section{Toward a CIFRE Doctoral Continuation}
\label{sec:conc-phd}
%------------------------------------------------------------------------------

This thesis establishes both the technical foundation and the scientific positioning for a CIFRE doctoral research program. Three open scientific questions, identified through the limitations of this thesis, constitute a coherent doctoral research agenda:

\begin{enumerate}
\item \textbf{MCP as a formal enterprise AI integration standard:} Can MCP be formally analyzed as a security protocol, and what properties (confidentiality, integrity, access control) can be formally proven about MCP-mediated RAG architectures? This question connects to the formal methods literature and is motivated by the residual prompt injection risk identified in Section~\ref{sec:disc-security}.

\item \textbf{Scalable multi-source, multi-modal enterprise RAG:} How does the RAG+MCP architecture perform when federated across multiple enterprise data sources (SharePoint, email, databases) and multiple modalities (text, tables, structured data)? This question is motivated by the single-source limitation of this thesis and by the multi-server extension capability in the MCP specification.

\item \textbf{Enterprise AI adoption factors:} What organizational, cognitive, and system design factors determine whether employees adopt and trust AI-powered knowledge assistants? What is the relationship between response latency, faithfulness, citation quality, and adoption intent over time? This question connects to information systems adoption literature (TAM, UTAUT) and is motivated by the time-limited user study of this thesis.
\end{enumerate}

A CIFRE PhD co-supervised between Aivancity and the DNSI's host organization would provide the industrial context, data access, and organizational scope required to address all three questions.

%------------------------------------------------------------------------------
\section{Closing Remarks}
\label{sec:conc-remarks}
%------------------------------------------------------------------------------

The AntiGravity system addresses a real operational problem: the gap between the organizational knowledge embedded in enterprise documentation and the practical ability of employees to access and use that knowledge. The RAG+MCP architecture presented in this thesis is one technically grounded response to that problem. Its evaluation demonstrates feasibility under real DNSI production constraints --- corporate proxy, OAuth 2.0 authentication, mixed-language corpus, CPU-only deployment --- and identifies the primary path to production readiness: GPU infrastructure for inference latency and real-time delta synchronization for corpus freshness.

The contributions of this thesis are bounded by the constraints of a Master's PFE: a single organization, a proof-of-concept deployment, and a time-limited evaluation. These boundaries are acknowledged explicitly in Chapter~\ref{chap:discussion}. Within these boundaries, the thesis engages the research questions with methodological rigor and produces genuinely replicable artifacts. That is what a PFE contribution requires.


%---- BACK MATTER -------------------------------------------------------------
\printbibliography[heading=bibintoc, title={Bibliography}]

\appendix
%==============================================================================
% APPENDIX A — BENCHMARK DATASET
%==============================================================================
\chapter{Benchmark Dataset: Question-Answer Pairs}
\label{app:benchmark}

This appendix documents the structure and annotation guidelines for the 50-item benchmark dataset used in the evaluation of Chapter~\ref{chap:evaluation}. The full dataset is provided in machine-readable format (JSON Lines) alongside this thesis.

\section*{Dataset Structure}

Each record in the dataset has the following structure:

\begin{lstlisting}[language=Python, caption={Benchmark dataset record schema}]
{
  "id": "Q001",
  "category": "factoid" | "multi-hop" | "procedural" | "unanswerable",
  "query": "Natural language question text",
  "ground_truth_answer": "Expected answer text",
  "relevant_chunk_ids": ["chunk_id_1", "chunk_id_2"],
  "source_document_titles": ["Document Title 1"],
  "annotator_1_label": 1,
  "annotator_2_label": 1,
  "cohen_kappa_pair": 0.85,
  "notes": "Optional annotation notes"
}
\end{lstlisting}

\section*{Annotation Guidelines}

Annotators were provided the following guidelines:

\begin{enumerate}
  \item Read the question and the candidate document chunk.
  \item Label the chunk as \textbf{1 (Relevant)} if it contains information necessary or sufficient to answer the question, or \textbf{0 (Not Relevant)} if it does not.
  \item A chunk may be labeled Relevant even if it provides only partial information (relevant to multi-hop synthesis).
  \item For unanswerable questions: all chunks should be labeled 0 if the question cannot be answered from the corpus.
  \item When in doubt, label 0 and note the ambiguity.
\end{enumerate}

\textit{Note: The DNSI-specific document content and individual Q\&A pairs are omitted from this appendix due to the confidentiality agreement. The anonymized synthetic test corpus is provided in the digital supplementary materials.}

%==============================================================================
% APPENDIX B — USER STUDY QUESTIONNAIRE
%==============================================================================
\chapter{User Study Instruments}
\label{app:user-study}

\section*{SUS Questionnaire (System Usability Scale)}

Participants rated the following 10 statements on a 5-point Likert scale (1 = Strongly Disagree, 5 = Strongly Agree):

\begin{enumerate}
  \item I think that I would like to use this system frequently.
  \item I found the system unnecessarily complex.
  \item I thought the system was easy to use.
  \item I think that I would need the support of a technical person to be able to use this system.
  \item I found the various functions in this system were well integrated.
  \item I thought there was too much inconsistency in this system.
  \item I would imagine that most people would learn to use this system very quickly.
  \item I found the system very cumbersome to use.
  \item I felt very confident using the system.
  \item I needed to learn a lot of things before I could get going with this system.
\end{enumerate}

\section*{TAM Questionnaire (Technology Acceptance Model)}

\textbf{Perceived Usefulness (PU)} --- 5 items, 7-point Likert (1 = Extremely Unlikely, 7 = Extremely Likely):
\begin{enumerate}
  \item Using AntiGravity improves my performance in finding information.
  \item Using AntiGravity increases my productivity.
  \item Using AntiGravity enhances my effectiveness on the job.
  \item AntiGravity is useful in my job.
  \item Using AntiGravity saves me time when searching for information.
\end{enumerate}

\textbf{Perceived Ease of Use (PEOU)} --- 4 items, 7-point Likert:
\begin{enumerate}
  \item Learning to operate AntiGravity is easy for me.
  \item I find it easy to get AntiGravity to do what I want it to do.
  \item My interaction with AntiGravity is clear and understandable.
  \item I find AntiGravity easy to use.
\end{enumerate}

\section*{Task Scenarios}

The five task scenarios used in the task-based evaluation:

\begin{enumerate}
  \item \textbf{T1 (Factoid):} ``Find the retention period for project closure reports in the DNSI document management policy.''
  \item \textbf{T2 (Multi-hop):} ``What are the security requirements for deploying a new application that processes personal data on the CERP infrastructure?''
  \item \textbf{T3 (Procedural):} ``Find the step-by-step procedure for requesting access to a new SharePoint site.''
  \item \textbf{T4 (Multi-hop):} ``Compare the roles and responsibilities of the Project Manager and the Technical Lead in the DNSI project governance framework.''
  \item \textbf{T5 (Unanswerable):} ``What is the direct phone number of the DNSI's external cybersecurity auditor?''
\end{enumerate}

\section*{Semi-Structured Interview Guide}

\begin{enumerate}
  \item How would you describe your overall experience using AntiGravity?
  \item For which types of tasks would you prefer to use AntiGravity over SharePoint search?
  \item What features did you find most valuable? Why?
  \item What aspects of the system frustrated you?
  \item How much do you trust the responses provided by AntiGravity? What factors influence your trust?
  \item Would you use AntiGravity in your daily work if it were permanently deployed? Under what conditions?
\end{enumerate}

%==============================================================================
% APPENDIX C — SYSTEM CONFIGURATION AND DEPLOYMENT
%==============================================================================
\chapter{System Configuration and Deployment Guide}
\label{app:deployment}

\section*{Environment Specification}

The complete Python dependency specification is available in \texttt{requirements.txt} in the project repository. Key pinned versions:

\begin{lstlisting}[caption={Key dependency versions}]
streamlit==1.28.2
sentence-transformers==2.2.2
faiss-cpu==1.7.4
mcp==1.0.0
msal==1.24.1
pdfplumber==0.10.3
python-docx==1.1.0
llama-cpp-python==0.2.38
httpx==0.25.2
\end{lstlisting}

\section*{Docker Compose Topology}

\begin{lstlisting}[language=yaml, caption={Docker Compose service configuration (simplified)}]
services:
  mcp-server:
    build: ./mcp_server
    environment:
      - FAISS_INDEX_PATH=/data/index
      - CHUNK_STORE_PATH=/data/chunks.jsonl
      - PERMISSION_MANIFEST=/data/permissions.json
    volumes:
      - ./data:/data
    ports:
      - "8080:8080"

  streamlit-app:
    build: ./streamlit_app
    environment:
      - MCP_SERVER_URL=http://mcp-server:8080
      - AZURE_CLIENT_ID=${AZURE_CLIENT_ID}
      - AZURE_TENANT_ID=${AZURE_TENANT_ID}
      - MISTRAL_MODEL_PATH=/models/mistral-7b-instruct-v0.2.Q4_K_M.gguf
    depends_on:
      - mcp-server
    ports:
      - "8501:8501"

  ingestion-worker:
    build: ./ingestion
    environment:
      - SHAREPOINT_SITE_URL=${SHAREPOINT_SITE_URL}
      - AZURE_CLIENT_SECRET=${AZURE_CLIENT_SECRET}
    volumes:
      - ./data:/data
\end{lstlisting}

\section*{Configuration Parameters}

\begin{table}[H]
\centering
\caption{Tunable system configuration parameters}
\begin{tabularx}{\textwidth}{@{}lXll@{}}
\toprule
\textbf{Parameter} & \textbf{Description} & \textbf{Default} & \textbf{Range} \\
\midrule
\texttt{CHUNK\_SIZE\_TOKENS} & Target chunk size in tokens & 512 & 256--1024 \\
\texttt{CHUNK\_OVERLAP\_TOKENS} & Overlap between adjacent chunks & 64 & 0--128 \\
\texttt{RETRIEVAL\_K} & Over-retrieval candidates before ACL filter & 20 & 10--50 \\
\texttt{RERANK\_TOP\_K} & Final context chunks after re-ranking & 5 & 3--10 \\
\texttt{LLM\_TEMPERATURE} & Generation temperature & 0.1 & 0.0--0.5 \\
\texttt{LLM\_MAX\_TOKENS} & Maximum generation length & 512 & 256--1024 \\
\texttt{PERMISSION\_CACHE\_TTL} & User permission cache TTL (seconds) & 300 & 60--3600 \\
\bottomrule
\end{tabularx}
\end{table}

\section*{Mistral 7B Model Provenance}

Model: \texttt{mistral-7b-instruct-v0.2.Q4\_K\_M.gguf}\\
Source: \texttt{TheBloke/Mistral-7B-Instruct-v0.2-GGUF} (Hugging Face)\\
SHA-256: \texttt{[to be populated from actual download]}\\
License: Apache 2.0

%==============================================================================
% APPENDIX_D — PLACEHOLDER
% TODO: Content to be written
%==============================================================================
\chapter{[Title to be defined — see thesis\_outline.md]}
\label{chap:appendix-d}

\textit{This chapter is under construction. See \texttt{PFE\_FINAL/02\_Research\_Design/thesis\_outline.md} for the planned content of this chapter.}


%==============================================================================
% APPENDIX E — DECLARATION OF AI TOOL USE
%==============================================================================
\chapter{Declaration of Generative AI Tool Use}
\label{app:ai-declaration}

In accordance with Aivancity School for Technology and Business's policy on the use of generative AI tools in academic work \citep{AivancityGuide2025}, the following declaration documents all uses of generative AI in the production of this thesis.

\section*{Tool Used}

\textbf{Claude} (claude-sonnet-4-6), developed by Anthropic. Accessed via the Claude Code CLI interface.

\section*{Scope of Use}

Generative AI assistance was used for the following purposes:

\begin{enumerate}
  \item \textbf{Structural drafting:} Initial drafts of thesis chapter structures, section headings, and outline frameworks were generated with AI assistance and subsequently revised, reorganized, and substantially expanded by the author.
  \item \textbf{Literature synthesis support:} AI assistance was used to suggest connections between cited works and to draft initial paragraph structures for the literature review. All factual claims, citations, and analytical judgments were verified and authored by the student.
  \item \textbf{LaTeX formatting:} Table structures, equation formatting, and bibliographic entry formatting were assisted by AI tools.
  \item \textbf{Language editing:} Grammar, sentence structure, and clarity improvements were suggested by AI tools and accepted or modified by the author.
\end{enumerate}

\section*{Scope of Non-Use}

The following elements were not generated by AI and represent the student's own intellectual work:

\begin{itemize}
  \item The research question formulation and hypothesis design
  \item All architectural design decisions and their justifications
  \item The benchmark dataset construction and annotation
  \item The experimental results and their interpretation
  \item The critical analysis of limitations and future work
  \item The code implementation of the AntiGravity system
  \item All final editorial decisions regarding content, emphasis, and argumentation
\end{itemize}

\section*{Compliance Statement}

This declaration complies with rules 1 (full transparency), 2 (systematic reference), 3 (fair use), and 4 (ethics) as specified in Chapter III of the Aivancity PFE Guidelines \citep{AivancityGuide2025}. The use of generative AI in this thesis is assistive and supplementary; the thesis primarily reflects the student's own skills, ideas, and efforts.

\bigskip

\noindent Signed: \underline{\hspace{5cm}} \hspace{2cm} Date: \underline{\hspace{3cm}}

\noindent Jean Direl NZE


\end{document}
